\worldchapter{\trans{templates_title}}{appendix-templates}
\label{ch:appendix-templates}

\begin{multicols}{2}

\section{Abstammung}

Die Charakterschablonen der Kategorie Abstammung beschreiben die Herkunft des Charakters. Bei der Charaktererschaffung kann eine Schablone dieser Kategorie kostenlos notiert werden. Die Schablonen dieser Kategorie können jedoch nicht im weiteren Verlauf der Charakterentwicklung gegen Reputation erworben werden.

\begin{tcolorbox}
    \subsection*{Mensch \hfill \small{0 \trans{reputation_label}}}
    \tcblower




\begin{tabular}{@{}ll@{}}
    \textbf{Bonuswürfel} & +2\\
    \textbf{Wiederholungswürfe} & +2\\
\end{tabular}

\wbif{phasesix}{
    \vspace{3mm}
\faIcon{cube}\hspace{0.5em}}{}

\end{tcolorbox}

\section{Beruf}

Die Charakterschablonen dieser Kategorie stellen Berufe dar, in denen der Charakter eine größere Professionalität erlangt hat. Die Schablonen sind in der Regel teurer als andere, bieten jedoch jeweils ein Paket aus Grundwissen und Fertigkeiten.

\begin{tcolorbox}
    \subsection*{Spion \hfill \small{11 \trans{reputation_label}}}
    \tcblower
\textit{Rivale}: Du hast einen Rivalen. Irgendjemand konkurriert mit dir, sei es auf einer beruflichen, professionellen oder tödlichen Ebene.


\vspace{3mm}

\begin{tabular}{@{}ll@{}}
    \textbf{Schnelligkeit} & +2\\
    \textbf{Wahrnehmung} & +2\\
    \textbf{Heimlichkeit} & +2\\
    \textbf{Kommunikation} & +1\\
    \textbf{Täuschen und Betrügen} & +2\\
\end{tabular}

\wbif{phasesix}{
    \vspace{3mm}
\faIcon{cube}\hspace{0.5em}}{}

\end{tcolorbox}

\begin{tcolorbox}
    \subsection*{Förster \hfill \small{7 \trans{reputation_label}}}
    \tcblower




\begin{tabular}{@{}ll@{}}
    \textbf{Kraft} & +1\\
    \textbf{Naturkunde} & +2\\
    \textbf{Orientierung} & +2\\
    \textbf{Kommunikation} & -1\\
    \textbf{Schießen} & +1\\
\end{tabular}

\wbif{phasesix}{
    \vspace{3mm}
\faIcon{cube}\hspace{0.5em}}{}

\end{tcolorbox}

\begin{tcolorbox}
    \subsection*{Banker \hfill \small{7 \trans{reputation_label}}}
    \tcblower
\textit{Gierig}: Immer wenn es darum geht den eigenen Reichtum zu mehren wirfst du auf deine Gewissenhaftigkeit. Misslingt der Wurf, wählst du den Weg des Geldes.


\vspace{3mm}

\begin{tabular}{@{}ll@{}}
    \textbf{Logik} & +2\\
    \textbf{Täuschen und Betrügen} & +1\\
    \textbf{Buchhaltung} & +2\\
\end{tabular}

\wbif{phasesix}{
    \vspace{3mm}
\faIcon{umbrella}\hspace{0.5em}\faIcon{flag}\hspace{0.5em}\faIcon{plane}\hspace{0.5em}\faIcon{microchip}\hspace{0.5em}}{}

\end{tcolorbox}

\begin{tcolorbox}
    \subsection*{Boxer \hfill \small{8 \trans{reputation_label}}}
    \tcblower




\begin{tabular}{@{}ll@{}}
    \textbf{Ausdauer} & +2\\
    \textbf{Schnelligkeit} & +2\\
    \textbf{Nahkampf} & +2\\
\end{tabular}

\wbif{phasesix}{
    \vspace{3mm}
\faIcon{cube}\hspace{0.5em}}{}

\end{tcolorbox}

\begin{tcolorbox}
    \subsection*{Lehrer \hfill \small{9 \trans{reputation_label}}}
    \tcblower




\begin{tabular}{@{}ll@{}}
    \textbf{Bildung} & +3\\
    \textbf{Auffassungsgabe} & +1\\
    \textbf{Kommunikation} & +2\\
    \textbf{Historie} & +1\\
\end{tabular}

\wbif{phasesix}{
    \vspace{3mm}
\faIcon{umbrella}\hspace{0.5em}\faIcon{flag}\hspace{0.5em}\faIcon{plane}\hspace{0.5em}\faIcon{microchip}\hspace{0.5em}}{}

\end{tcolorbox}

\begin{tcolorbox}
    \subsection*{Soldat \hfill \small{14 \trans{reputation_label}}}
    \tcblower
\textit{Obrigkeitshörig}: Du befolgst jeden Befehl deines Vorgesetzten, ohne weiter darüber nachzudenken.


\vspace{3mm}

\begin{tabular}{@{}ll@{}}
    \textbf{Ausdauer} & +1\\
    \textbf{Kraft} & +1\\
    \textbf{Orientierung} & +1\\
    \textbf{Erste Hilfe} & +1\\
    \textbf{Tapferkeit} & +1\\
    \textbf{Fahren} & +1\\
    \textbf{Werfen} & +2\\
    \textbf{Schießen} & +2\\
    \textbf{Nahkampf} & +1\\
    \textbf{Fahrzeuge} & +1\\
\end{tabular}

\wbif{phasesix}{
    \vspace{3mm}
\faIcon{cube}\hspace{0.5em}}{}

\end{tcolorbox}

\begin{tcolorbox}
    \subsection*{Geistlicher \hfill \small{10 \trans{reputation_label}}}
    \tcblower
\textit{Religiös}: Du bist religiös, glaubst an deine Gottheit und verteidigst deinen Glauben auch aktiv.


\vspace{3mm}

\begin{tabular}{@{}ll@{}}
    \textbf{Schicksalswürfel} & +2\\
    \textbf{Bildung} & +1\\
    \textbf{Religion} & +3\\
    \textbf{Kommunikation} & +2\\
\end{tabular}

\wbif{phasesix}{
    \vspace{3mm}
\faIcon{cube}\hspace{0.5em}}{}

\end{tcolorbox}

\begin{tcolorbox}
    \subsection*{Schmied \hfill \small{8 \trans{reputation_label}}}
    \tcblower




\begin{tabular}{@{}ll@{}}
    \textbf{Kraft} & +2\\
    \textbf{Geschick} & +2\\
    \textbf{Mechanik} & +2\\
\end{tabular}

\wbif{phasesix}{
    \vspace{3mm}
\faIcon{chess-rook}\hspace{0.5em}\faIcon{umbrella}\hspace{0.5em}\faIcon{flag}\hspace{0.5em}\faIcon{plane}\hspace{0.5em}}{}

\end{tcolorbox}

\begin{tcolorbox}
    \subsection*{Schreiber \hfill \small{8 \trans{reputation_label}}}
    \tcblower




\begin{tabular}{@{}ll@{}}
    \textbf{Gewissenhaftigkeit} & +1\\
    \textbf{Bildung} & +1\\
    \textbf{Kommunikation} & +2\\
    \textbf{Lesen/Schreiben} & +2\\
\end{tabular}

\wbif{phasesix}{
    \vspace{3mm}
\faIcon{chess-rook}\hspace{0.5em}}{}

\end{tcolorbox}

\begin{tcolorbox}
    \subsection*{Söldner \hfill \small{8 \trans{reputation_label}}}
    \tcblower




\begin{tabular}{@{}ll@{}}
    \textbf{Ausdauer} & +1\\
    \textbf{Geschick} & +1\\
    \textbf{Überzeugen} & +1\\
    \textbf{Schießen} & +1\\
    \textbf{Nahkampf} & +2\\
\end{tabular}

\wbif{phasesix}{
    \vspace{3mm}
\faIcon{cube}\hspace{0.5em}}{}

\end{tcolorbox}

\begin{tcolorbox}
    \subsection*{Informatiker \hfill \small{11 \trans{reputation_label}}}
    \tcblower




\begin{tabular}{@{}ll@{}}
    \textbf{Gewissenhaftigkeit} & +2\\
    \textbf{Logik} & +3\\
    \textbf{Mechanik} & +1\\
    \textbf{Informatik} & +3\\
\end{tabular}

\wbif{phasesix}{
    \vspace{3mm}
\faIcon{plane}\hspace{0.5em}\faIcon{microchip}\hspace{0.5em}}{}

\end{tcolorbox}

\begin{tcolorbox}
    \subsection*{Sozialpädagoge \hfill \small{9 \trans{reputation_label}}}
    \tcblower
\textit{Gutmensch}: Der Charakter will nur gutes tun, anderen helfen und ist generell selbstlos.


\vspace{3mm}

\begin{tabular}{@{}ll@{}}
    \textbf{Ausdauer} & +2\\
    \textbf{Charme} & +1\\
    \textbf{Bildung} & +2\\
    \textbf{Kommunikation} & +2\\
\end{tabular}

\wbif{phasesix}{
    \vspace{3mm}
\faIcon{plane}\hspace{0.5em}\faIcon{microchip}\hspace{0.5em}}{}

\end{tcolorbox}

\begin{tcolorbox}
    \subsection*{Frührentner \hfill \small{6 \trans{reputation_label}}}
    \tcblower
\textit{Wutbürger}: Du hast eine aggresive Grundhaltung gegenüber allem. Du schreibst Falschparker auf, motzt über die Regierung und postest Unsinn auf Boomerbook.


\vspace{3mm}

\begin{tabular}{@{}ll@{}}
    \textbf{Resistenz} & -1\\
    \textbf{Willenskraft} & +1\\
    \textbf{Einschüchtern} & +2\\
    \textbf{Recht} & +2\\
\end{tabular}

\wbif{phasesix}{
    \vspace{3mm}
\faIcon{microchip}\hspace{0.5em}}{}

\end{tcolorbox}

\begin{tcolorbox}
    \subsection*{Seefahrer \hfill \small{12 \trans{reputation_label}}}
    \tcblower




\begin{tabular}{@{}ll@{}}
    \textbf{Geschick} & +1\\
    \textbf{Kraft} & +1\\
    \textbf{Resistenz} & +1\\
    \textbf{Fahren} & +1\\
    \textbf{Nahkampf} & +1\\
    \textbf{Werfen} & +2\\
    \textbf{Seefahrt} & +3\\
\end{tabular}

\wbif{phasesix}{
    \vspace{3mm}
\faIcon{cube}\hspace{0.5em}}{}

\end{tcolorbox}

\begin{tcolorbox}
    \subsection*{Kumpel \hfill \small{10 \trans{reputation_label}}}
    \tcblower




\begin{tabular}{@{}ll@{}}
    \textbf{Kraft} & +2\\
    \textbf{Ausdauer} & +1\\
    \textbf{Tapferkeit} & +1\\
    \textbf{Sprengstoffe} & +2\\
    \textbf{Gesteinskunde} & +2\\
\end{tabular}

\wbif{phasesix}{
    \vspace{3mm}
\faIcon{umbrella}\hspace{0.5em}\faIcon{flag}\hspace{0.5em}\faIcon{plane}\hspace{0.5em}\faIcon{microchip}\hspace{0.5em}}{}

\end{tcolorbox}

\begin{tcolorbox}
    \subsection*{Kurier \hfill \small{8 \trans{reputation_label}}}
    \tcblower




\begin{tabular}{@{}ll@{}}
    \textbf{Schnelligkeit} & +2\\
    \textbf{Ausdauer} & +1\\
    \textbf{Fahren} & +1\\
    \textbf{Orientierung} & +2\\
\end{tabular}

\wbif{phasesix}{
    \vspace{3mm}
\faIcon{cube}\hspace{0.5em}}{}

\end{tcolorbox}

\begin{tcolorbox}
    \subsection*{Händler \hfill \small{8 \trans{reputation_label}}}
    \tcblower




\begin{tabular}{@{}ll@{}}
    \textbf{Logik} & +2\\
    \textbf{Auffassungsgabe} & +1\\
    \textbf{Empathie} & +1\\
    \textbf{Überzeugen} & +2\\
\end{tabular}

\wbif{phasesix}{
    \vspace{3mm}
\faIcon{cube}\hspace{0.5em}}{}

\end{tcolorbox}

\begin{tcolorbox}
    \subsection*{Mediengestalter \hfill \small{5 \trans{reputation_label}}}
    \tcblower




\begin{tabular}{@{}ll@{}}
    \textbf{Kraft} & -1\\
    \textbf{Auffassungsgabe} & +1\\
    \textbf{Untersuchen} & +2\\
    \textbf{Wahrnehmung} & +1\\
\end{tabular}

\wbif{phasesix}{
    \vspace{3mm}
\faIcon{plane}\hspace{0.5em}\faIcon{microchip}\hspace{0.5em}}{}

\end{tcolorbox}

\begin{tcolorbox}
    \subsection*{Mechatroniker \hfill \small{9 \trans{reputation_label}}}
    \tcblower




\begin{tabular}{@{}ll@{}}
    \textbf{Geschick} & +2\\
    \textbf{Mechanik} & +2\\
    \textbf{Fahren} & +1\\
    \textbf{Fahrzeuge} & +2\\
\end{tabular}

\wbif{phasesix}{
    \vspace{3mm}
\faIcon{plane}\hspace{0.5em}\faIcon{microchip}\hspace{0.5em}}{}

\end{tcolorbox}

\begin{tcolorbox}
    \subsection*{Landwirt \hfill \small{7 \trans{reputation_label}}}
    \tcblower




\begin{tabular}{@{}ll@{}}
    \textbf{Gewissenhaftigkeit} & -1\\
    \textbf{Kommunikation} & -1\\
    \textbf{Fahren} & +2\\
\end{tabular}

\wbif{phasesix}{
    \vspace{3mm}
\faIcon{umbrella}\hspace{0.5em}\faIcon{flag}\hspace{0.5em}\faIcon{plane}\hspace{0.5em}\faIcon{microchip}\hspace{0.5em}}{}

\end{tcolorbox}

\begin{tcolorbox}
    \subsection*{Friseur \hfill \small{6 \trans{reputation_label}}}
    \tcblower




\begin{tabular}{@{}ll@{}}
    \textbf{Geschick} & +1\\
    \textbf{Charme} & +2\\
    \textbf{Empathie} & +1\\
\end{tabular}

\wbif{phasesix}{
    \vspace{3mm}
\faIcon{umbrella}\hspace{0.5em}\faIcon{flag}\hspace{0.5em}\faIcon{plane}\hspace{0.5em}\faIcon{microchip}\hspace{0.5em}}{}

\end{tcolorbox}

\begin{tcolorbox}
    \subsection*{Entertainer \hfill \small{11 \trans{reputation_label}}}
    \tcblower




\begin{tabular}{@{}ll@{}}
    \textbf{Schicksalswürfel} & +1\\
    \textbf{Charme} & +2\\
    \textbf{Kommunikation} & +1\\
    \textbf{Nahkampf} & +1\\
    \textbf{Darbietung} & +2\\
    \textbf{Täuschen und Betrügen} & +2\\
\end{tabular}

\wbif{phasesix}{
    \vspace{3mm}
\faIcon{umbrella}\hspace{0.5em}\faIcon{flag}\hspace{0.5em}\faIcon{plane}\hspace{0.5em}\faIcon{microchip}\hspace{0.5em}}{}

\end{tcolorbox}

\begin{tcolorbox}
    \subsection*{Journalist \hfill \small{8 \trans{reputation_label}}}
    \tcblower




\begin{tabular}{@{}ll@{}}
    \textbf{Schicksalswürfel} & +1\\
    \textbf{Bildung} & +1\\
    \textbf{Untersuchen} & +2\\
    \textbf{Kommunikation} & +2\\
\end{tabular}

\wbif{phasesix}{
    \vspace{3mm}
\faIcon{umbrella}\hspace{0.5em}\faIcon{flag}\hspace{0.5em}\faIcon{plane}\hspace{0.5em}\faIcon{microchip}\hspace{0.5em}\faIcon{space-shuttle}\hspace{0.5em}}{}

\end{tcolorbox}

\begin{tcolorbox}
    \subsection*{Polizist \hfill \small{8 \trans{reputation_label}}}
    \tcblower
\textit{Obrigkeitshörig}: Du befolgst jeden Befehl deines Vorgesetzten, ohne weiter darüber nachzudenken.


\vspace{3mm}

\begin{tabular}{@{}ll@{}}
    \textbf{Ausdauer} & +1\\
    \textbf{Gewissenhaftigkeit} & +1\\
    \textbf{Überzeugen} & +1\\
    \textbf{Kommunikation} & +1\\
    \textbf{Schießen} & +2\\
\end{tabular}

\wbif{phasesix}{
    \vspace{3mm}
\faIcon{umbrella}\hspace{0.5em}\faIcon{flag}\hspace{0.5em}\faIcon{plane}\hspace{0.5em}\faIcon{microchip}\hspace{0.5em}\faIcon{space-shuttle}\hspace{0.5em}}{}

\end{tcolorbox}

\begin{tcolorbox}
    \subsection*{Gangmitglied \hfill \small{5 \trans{reputation_label}}}
    \tcblower
\textit{Rivale}: Du hast einen Rivalen. Irgendjemand konkurriert mit dir, sei es auf einer beruflichen, professionellen oder tödlichen Ebene.


\vspace{3mm}

\begin{tabular}{@{}ll@{}}
    \textbf{Resistenz} & +1\\
    \textbf{Bildung} & -1\\
    \textbf{Nahkampf} & +2\\
    \textbf{Schießen} & +1\\
\end{tabular}

\wbif{phasesix}{
    \vspace{3mm}
\faIcon{umbrella}\hspace{0.5em}\faIcon{flag}\hspace{0.5em}\faIcon{plane}\hspace{0.5em}\faIcon{microchip}\hspace{0.5em}\faIcon{space-shuttle}\hspace{0.5em}}{}

\end{tcolorbox}

\begin{tcolorbox}
    \subsection*{Sanitäter \hfill \small{8 \trans{reputation_label}}}
    \tcblower




\begin{tabular}{@{}ll@{}}
    \textbf{Maximaler Stress} & +1\\
    \textbf{Gewissenhaftigkeit} & +1\\
    \textbf{Erste Hilfe} & +3\\
    \textbf{Medizin} & +1\\
\end{tabular}

\wbif{phasesix}{
    \vspace{3mm}
\faIcon{umbrella}\hspace{0.5em}\faIcon{flag}\hspace{0.5em}\faIcon{plane}\hspace{0.5em}\faIcon{microchip}\hspace{0.5em}}{}

\end{tcolorbox}

\begin{tcolorbox}
    \subsection*{Fahrer \hfill \small{8 \trans{reputation_label}}}
    \tcblower




\begin{tabular}{@{}ll@{}}
    \textbf{Fahren} & +4\\
    \textbf{Fahrzeuge} & +2\\
\end{tabular}

\wbif{phasesix}{
    \vspace{3mm}
\faIcon{umbrella}\hspace{0.5em}\faIcon{flag}\hspace{0.5em}\faIcon{plane}\hspace{0.5em}\faIcon{microchip}\hspace{0.5em}}{}

\end{tcolorbox}

\begin{tcolorbox}
    \subsection*{Bestatter \hfill \small{8 \trans{reputation_label}}}
    \tcblower




\begin{tabular}{@{}ll@{}}
    \textbf{Ausdauer} & +1\\
    \textbf{Charme} & +1\\
    \textbf{Fahren} & +1\\
    \textbf{Empathie} & +1\\
    \textbf{Etikette} & +1\\
\end{tabular}

\wbif{phasesix}{
    \vspace{3mm}
\faIcon{cube}\hspace{0.5em}}{}

\end{tcolorbox}

\begin{tcolorbox}
    \subsection*{Gelehrter \hfill \small{10 \trans{reputation_label}}}
    \tcblower




\begin{tabular}{@{}ll@{}}
    \textbf{Bildung} & +4\\
    \textbf{Naturkunde} & +1\\
    \textbf{Historie} & +2\\
    \textbf{Kommunikation} & +1\\
\end{tabular}

\wbif{phasesix}{
    \vspace{3mm}
\faIcon{chess-rook}\hspace{0.5em}}{}

\end{tcolorbox}

\begin{tcolorbox}
    \subsection*{Abenteurer \hfill \small{9 \trans{reputation_label}}}
    \tcblower




\begin{tabular}{@{}ll@{}}
    \textbf{Wiederholungswürfe} & +1\\
    \textbf{Kraft} & +1\\
    \textbf{Geschick} & +1\\
    \textbf{Orientierung} & +2\\
    \textbf{Nahkampf} & +1\\
    \textbf{Untersuchen} & +1\\
\end{tabular}

\wbif{phasesix}{
    \vspace{3mm}
\faIcon{chess-rook}\hspace{0.5em}\faIcon{umbrella}\hspace{0.5em}}{}

\end{tcolorbox}

\begin{tcolorbox}
    \subsection*{Barde \hfill \small{12 \trans{reputation_label}}}
    \tcblower




\begin{tabular}{@{}ll@{}}
    \textbf{Charme} & +2\\
    \textbf{Attraktivität} & +2\\
    \textbf{Historie} & +2\\
    \textbf{Darbietung} & +3\\
    \textbf{Lesen/Schreiben} & +1\\
\end{tabular}

\wbif{phasesix}{
    \vspace{3mm}
\faIcon{chess-rook}\hspace{0.5em}}{}

\end{tcolorbox}

\begin{tcolorbox}
    \subsection*{Dieb \hfill \small{10 \trans{reputation_label}}}
    \tcblower




\begin{tabular}{@{}ll@{}}
    \textbf{Geschick} & +2\\
    \textbf{Heimlichkeit} & +3\\
    \textbf{Akrobatik} & +1\\
    \textbf{Nahkampf} & +2\\
\end{tabular}

\wbif{phasesix}{
    \vspace{3mm}
\faIcon{cube}\hspace{0.5em}}{}

\end{tcolorbox}

\begin{tcolorbox}
    \subsection*{Meuchler \hfill \small{11 \trans{reputation_label}}}
    \tcblower




\begin{tabular}{@{}ll@{}}
    \textbf{Geschick} & +1\\
    \textbf{Akrobatik} & +2\\
    \textbf{Schießen} & +2\\
    \textbf{Nahkampf} & +2\\
    \textbf{Täuschen und Betrügen} & +2\\
\end{tabular}

\wbif{phasesix}{
    \vspace{3mm}
\faIcon{cube}\hspace{0.5em}}{}

\end{tcolorbox}

\begin{tcolorbox}
    \subsection*{Medicus \hfill \small{10 \trans{reputation_label}}}
    \tcblower




\begin{tabular}{@{}ll@{}}
    \textbf{Geschick} & +1\\
    \textbf{Gewissenhaftigkeit} & +1\\
    \textbf{Erste Hilfe} & +4\\
    \textbf{Medizin} & +2\\
\end{tabular}

\wbif{phasesix}{
    \vspace{3mm}
\faIcon{chess-rook}\hspace{0.5em}}{}

\end{tcolorbox}

\begin{tcolorbox}
    \subsection*{Feinmechaniker \hfill \small{10 \trans{reputation_label}}}
    \tcblower




\begin{tabular}{@{}ll@{}}
    \textbf{Geschick} & +3\\
    \textbf{Gewissenhaftigkeit} & +2\\
    \textbf{Mechanik} & +3\\
\end{tabular}

\wbif{phasesix}{
    \vspace{3mm}
\faIcon{umbrella}\hspace{0.5em}\faIcon{flag}\hspace{0.5em}\faIcon{plane}\hspace{0.5em}\faIcon{microchip}\hspace{0.5em}}{}

\end{tcolorbox}

\begin{tcolorbox}
    \subsection*{Pilot \hfill \small{8 \trans{reputation_label}}}
    \tcblower




\begin{tabular}{@{}ll@{}}
    \textbf{Charme} & +1\\
    \textbf{Ausdauer} & +1\\
    \textbf{Auffassungsgabe} & +1\\
    \textbf{Tapferkeit} & +1\\
    \textbf{Luftfahrt} & +2\\
\end{tabular}

\wbif{phasesix}{
    \vspace{3mm}
\faIcon{flag}\hspace{0.5em}\faIcon{plane}\hspace{0.5em}\faIcon{microchip}\hspace{0.5em}\faIcon{space-shuttle}\hspace{0.5em}}{}

\end{tcolorbox}

\begin{tcolorbox}
    \subsection*{Ritter \hfill \small{10 \trans{reputation_label}}}
    \tcblower




\begin{tabular}{@{}ll@{}}
    \textbf{Kraft} & +1\\
    \textbf{Ausdauer} & +2\\
    \textbf{Geschick} & +1\\
    \textbf{Nahkampf} & +2\\
    \textbf{Politik} & +1\\
    \textbf{Reiten} & +1\\
\end{tabular}

\wbif{phasesix}{
    \vspace{3mm}
\faIcon{chess-rook}\hspace{0.5em}}{}

\end{tcolorbox}

\begin{tcolorbox}
    \subsection*{Politiker \hfill \small{7 \trans{reputation_label}}}
    \tcblower
\textit{Rivale}: Du hast einen Rivalen. Irgendjemand konkurriert mit dir, sei es auf einer beruflichen, professionellen oder tödlichen Ebene.

\begin{pquote}{Philipp Amthor}
Ich bin nicht käuflich. Gleichwohl habe ich mich politisch angreifbar gemacht und kann die Kritik nachvollziehen. Es war ein Fehler.
\end{pquote}

\vspace{3mm}

\begin{tabular}{@{}ll@{}}
    \textbf{Einschüchtern} & +1\\
    \textbf{Kommunikation} & +2\\
    \textbf{Politik} & +3\\
    \textbf{Täuschen und Betrügen} & +1\\
\end{tabular}

\wbif{phasesix}{
    \vspace{3mm}
\faIcon{umbrella}\hspace{0.5em}\faIcon{flag}\hspace{0.5em}\faIcon{plane}\hspace{0.5em}\faIcon{microchip}\hspace{0.5em}}{}

\end{tcolorbox}

\begin{tcolorbox}
    \subsection*{Programmierer \hfill \small{10 \trans{reputation_label}}}
    \tcblower


\begin{pquote}{Robert C. Martin}
See you are my tribe.

I don't care if you're young or you're old, or black or white, or a man or a woman. i don't care who you like or who you love. If you are a programmer, you are part of my tribe. You and I, we all together, share a passion for something. And we can communicate about it. In a way most other people can't. And so that's something we should cherish together.
\end{pquote}

\vspace{3mm}

\begin{tabular}{@{}ll@{}}
    \textbf{Logik} & +3\\
    \textbf{Hacken} & +2\\
    \textbf{Informatik} & +3\\
\end{tabular}

\wbif{phasesix}{
    \vspace{3mm}
\faIcon{plane}\hspace{0.5em}\faIcon{microchip}\hspace{0.5em}}{}

\end{tcolorbox}

\begin{tcolorbox}
    \subsection*{Bader \hfill \small{7 \trans{reputation_label}}}
    \tcblower




\begin{tabular}{@{}ll@{}}
    \textbf{Geschick} & +1\\
    \textbf{Erste Hilfe} & +2\\
    \textbf{Täuschen und Betrügen} & +1\\
    \textbf{Medizin} & +1\\
\end{tabular}

\wbif{phasesix}{
    \vspace{3mm}
\faIcon{chess-rook}\hspace{0.5em}}{}

\end{tcolorbox}

\begin{tcolorbox}
    \subsection*{Raumpirat \hfill \small{13 \trans{reputation_label}}}
    \tcblower
\textit{Gierig}: Immer wenn es darum geht den eigenen Reichtum zu mehren wirfst du auf deine Gewissenhaftigkeit. Misslingt der Wurf, wählst du den Weg des Geldes.


\vspace{3mm}

\begin{tabular}{@{}ll@{}}
    \textbf{Bonuswürfel} & +2\\
    \textbf{Auffassungsgabe} & +2\\
    \textbf{Charme} & +2\\
    \textbf{Einschüchtern} & +2\\
    \textbf{Schießen} & +1\\
    \textbf{Täuschen und Betrügen} & +1\\
    \textbf{Fahrzeuge} & +1\\
\end{tabular}

\wbif{phasesix}{
    \vspace{3mm}
\faIcon{space-shuttle}\hspace{0.5em}}{}

\end{tcolorbox}

\begin{tcolorbox}
    \subsection*{Bürokraft \hfill \small{9 \trans{reputation_label}}}
    \tcblower




\begin{tabular}{@{}ll@{}}
    \textbf{Auffassungsgabe} & +1\\
    \textbf{Gewissenhaftigkeit} & +1\\
    \textbf{Heimlichkeit} & +1\\
    \textbf{Überzeugen} & +2\\
    \textbf{Verwaltung} & +2\\
\end{tabular}

\wbif{phasesix}{
    \vspace{3mm}
\faIcon{flag}\hspace{0.5em}\faIcon{plane}\hspace{0.5em}\faIcon{microchip}\hspace{0.5em}\faIcon{space-shuttle}\hspace{0.5em}}{}

\end{tcolorbox}

\begin{tcolorbox}
    \subsection*{Abdecker \hfill \small{5 \trans{reputation_label}}}
    \tcblower




\begin{tabular}{@{}ll@{}}
    \textbf{Resistenz} & +1\\
    \textbf{Attraktivität} & -1\\
    \textbf{Einschüchtern} & +1\\
    \textbf{Naturkunde} & +1\\
    \textbf{Medizin} & +1\\
\end{tabular}

\wbif{phasesix}{
    \vspace{3mm}
\faIcon{eye}\hspace{0.5em}\faIcon{chess-rook}\hspace{0.5em}\faIcon{umbrella}\hspace{0.5em}}{}

\end{tcolorbox}

\begin{tcolorbox}
    \subsection*{Geisterjäger \hfill \small{8 \trans{reputation_label}}}
    \tcblower


\begin{pquote}{Sam Winchester}
People don’t just disappear, Dean. Other people just stop looking for them.
\end{pquote}

\vspace{3mm}

\begin{tabular}{@{}ll@{}}
    \textbf{Maximaler Stress} & +2\\
    \textbf{Resistenz} & +1\\
    \textbf{Tapferkeit} & +1\\
    \textbf{Nahkampf} & +1\\
    \textbf{Altertümer} & +1\\
\end{tabular}

\wbif{phasesix}{
    \vspace{3mm}
\faIcon{spider}\hspace{0.5em}}{}

\end{tcolorbox}

\begin{tcolorbox}
    \subsection*{Wirt \hfill \small{9 \trans{reputation_label}}}
    \tcblower




\begin{tabular}{@{}ll@{}}
    \textbf{Auffassungsgabe} & +2\\
    \textbf{Empathie} & +2\\
    \textbf{Kommunikation} & +2\\
\end{tabular}

\wbif{phasesix}{
    \vspace{3mm}
\faIcon{cube}\hspace{0.5em}}{}

\end{tcolorbox}

\begin{tcolorbox}
    \subsection*{Archäologe \hfill \small{10 \trans{reputation_label}}}
    \tcblower




\begin{tabular}{@{}ll@{}}
    \textbf{Schicksalswürfel} & +1\\
    \textbf{Geschick} & +1\\
    \textbf{Wahrnehmung} & +2\\
    \textbf{Historie} & +1\\
    \textbf{Altertümer} & +4\\
\end{tabular}

\wbif{phasesix}{
    \vspace{3mm}
\faIcon{flag}\hspace{0.5em}\faIcon{plane}\hspace{0.5em}\faIcon{microchip}\hspace{0.5em}}{}

\end{tcolorbox}

\begin{tcolorbox}
    \subsection*{Autor \hfill \small{6 \trans{reputation_label}}}
    \tcblower


\begin{pquote}{Stephen King}
The road to hell is paved with adverbs.
\end{pquote}

\vspace{3mm}

\begin{tabular}{@{}ll@{}}
    \textbf{Bonuswürfel} & +1\\
    \textbf{Bildung} & +1\\
    \textbf{Gewissenhaftigkeit} & +1\\
    \textbf{Historie} & +1\\
\end{tabular}

\wbif{phasesix}{
    \vspace{3mm}
\faIcon{umbrella}\hspace{0.5em}\faIcon{flag}\hspace{0.5em}\faIcon{plane}\hspace{0.5em}\faIcon{microchip}\hspace{0.5em}\faIcon{space-shuttle}\hspace{0.5em}}{}

\end{tcolorbox}

\begin{tcolorbox}
    \subsection*{Fachkraft für Veranstaltungstechnik \hfill \small{10 \trans{reputation_label}}}
    \tcblower


\begin{pquote}{Phillip Schröder}
"Heute ist Open-End-Feierabend."
\end{pquote}

\vspace{3mm}

\begin{tabular}{@{}ll@{}}
    \textbf{Maximaler Stress} & +2\\
    \textbf{Gewissenhaftigkeit} & -1\\
    \textbf{Tapferkeit} & +1\\
    \textbf{Fahren} & +2\\
    \textbf{Recht} & +1\\
\end{tabular}

\wbif{phasesix}{
    \vspace{3mm}
\faIcon{microchip}\hspace{0.5em}}{}

\end{tcolorbox}

\begin{tcolorbox}
    \subsection*{Adel \hfill \small{10 \trans{reputation_label}}}
    \tcblower
\textit{Eitelkeit}: Du bist über alle Maßen eitel und zeigst dies auch oft und gerne.

\begin{pquote}{Friedrich Schiller}
Adel ist auch in der sittlichen Welt. Gemeine Naturen zahlen mit dem, was sie tun, edle mit dem, was sie sind.
\end{pquote}

\vspace{3mm}

\begin{tabular}{@{}ll@{}}
    \textbf{Charme} & +2\\
    \textbf{Überzeugen} & +2\\
    \textbf{Schießen} & +1\\
    \textbf{Empathie} & -1\\
    \textbf{Etikette} & +2\\
    \textbf{Reiten} & +1\\
\end{tabular}

\wbif{phasesix}{
    \vspace{3mm}
\faIcon{microchip}\hspace{0.5em}}{}

\end{tcolorbox}

\begin{tcolorbox}
    \subsection*{Bereiter \hfill \small{10 \trans{reputation_label}}}
    \tcblower




\begin{tabular}{@{}ll@{}}
    \textbf{Geschick} & +2\\
    \textbf{Kraft} & +1\\
    \textbf{Empathie} & +1\\
    \textbf{Reiten} & +3\\
    \textbf{Zoologie} & +2\\
\end{tabular}

\wbif{phasesix}{
    \vspace{3mm}
\faIcon{eye}\hspace{0.5em}\faIcon{chess-rook}\hspace{0.5em}\faIcon{umbrella}\hspace{0.5em}\faIcon{flag}\hspace{0.5em}}{}

\end{tcolorbox}

\begin{tcolorbox}
    \subsection*{Druide \hfill \small{10 \trans{reputation_label}}}
    \tcblower


\begin{pquote}{Robert Louis Stevenson}
It is not so much for its beauty that the forest makes a claim upon men’s hearts, as for that subtle something, that quality of air that emanation from old trees, that so wonderfully changes and renews a weary spirit.
\end{pquote}

\vspace{3mm}

\begin{tabular}{@{}ll@{}}
    \textbf{Max Arkana} & +3\\
    \textbf{Zauberpunkte} & +5\\
    \textbf{Resistenz} & +1\\
    \textbf{Willenskraft} & +1\\
    \textbf{Naturkunde} & +2\\
    \textbf{Zaubern} & +1\\
    \textbf{Altertümer} & +1\\
    \textbf{Elementarmagie} & True\\
\end{tabular}

\wbif{phasesix}{
    \vspace{3mm}
\faIcon{magic}\hspace{0.5em}}{}

\end{tcolorbox}

\begin{tcolorbox}
    \subsection*{Hexe \hfill \small{10 \trans{reputation_label}}}
    \tcblower


\begin{pquote}{William Shakespeare}
Eye of newt, and toe of frog,
Wool of bat, and tongue of dog,
Adder's fork, and blind-worm's sting,
Lizard's leg, and owlet's wing,—
For a charm of powerful trouble,
Like a hell-broth boil and bubble.
Double, double toil and trouble;
Fire burn, and caldron bubble.
\end{pquote}

\vspace{3mm}

\begin{tabular}{@{}ll@{}}
    \textbf{Max Arkana} & +3\\
    \textbf{Zauberpunkte} & +5\\
    \textbf{Charme} & +1\\
    \textbf{Auffassungsgabe} & +1\\
    \textbf{Magiekunde} & +1\\
    \textbf{Zaubern} & +1\\
    \textbf{Erste Hilfe} & +2\\
    \textbf{Verfluchungen} & True\\
\end{tabular}

\wbif{phasesix}{
    \vspace{3mm}
\faIcon{magic}\hspace{0.5em}}{}

\end{tcolorbox}

\begin{tcolorbox}
    \subsection*{Dämonologe \hfill \small{10 \trans{reputation_label}}}
    \tcblower


\begin{pquote}{Brenna Yovanoff}
The treachery of demons is nothing compared to the betrayal of an angel.
\end{pquote}

\vspace{3mm}

\begin{tabular}{@{}ll@{}}
    \textbf{Zauberpunkte} & +5\\
    \textbf{Max Arkana} & +3\\
    \textbf{Charme} & +1\\
    \textbf{Ausdauer} & +1\\
    \textbf{Magiekunde} & +1\\
    \textbf{Darbietung} & +1\\
    \textbf{Zaubern} & +1\\
    \textbf{Altertümer} & +1\\
    \textbf{Dämonologie} & True\\
\end{tabular}

\wbif{phasesix}{
    \vspace{3mm}
\faIcon{magic}\hspace{0.5em}}{}

\end{tcolorbox}

\begin{tcolorbox}
    \subsection*{Waldläufer \hfill \small{10 \trans{reputation_label}}}
    \tcblower




\begin{tabular}{@{}ll@{}}
    \textbf{Max Arkana} & +3\\
    \textbf{Zauberpunkte} & +5\\
    \textbf{Ausdauer} & +1\\
    \textbf{Schnelligkeit} & +1\\
    \textbf{Naturkunde} & +1\\
    \textbf{Magiekunde} & +1\\
    \textbf{Zaubern} & +1\\
    \textbf{Schamanismus} & True\\
\end{tabular}

\wbif{phasesix}{
    \vspace{3mm}
\faIcon{magic}\hspace{0.5em}}{}

\end{tcolorbox}

\begin{tcolorbox}
    \subsection*{Zauberer \hfill \small{10 \trans{reputation_label}}}
    \tcblower




\begin{tabular}{@{}ll@{}}
    \textbf{Zauberpunkte} & +5\\
    \textbf{Max Arkana} & +3\\
    \textbf{Logik} & +1\\
    \textbf{Auffassungsgabe} & +1\\
    \textbf{Zaubern} & +1\\
    \textbf{Magiekunde} & +1\\
    \textbf{Mathematik} & +1\\
    \textbf{Zauberei} & True\\
\end{tabular}

\wbif{phasesix}{
    \vspace{3mm}
\faIcon{magic}\hspace{0.5em}}{}

\end{tcolorbox}

\begin{tcolorbox}
    \subsection*{Nekrologe \hfill \small{10 \trans{reputation_label}}}
    \tcblower




\begin{tabular}{@{}ll@{}}
    \textbf{Max Arkana} & +3\\
    \textbf{Zauberpunkte} & +5\\
    \textbf{Willenskraft} & +1\\
    \textbf{Resistenz} & +1\\
    \textbf{Täuschen und Betrügen} & +1\\
    \textbf{Heimlichkeit} & +1\\
    \textbf{Altertümer} & +1\\
    \textbf{Nekrologie} & True\\
\end{tabular}

\wbif{phasesix}{
    \vspace{3mm}
\faIcon{magic}\hspace{0.5em}}{}

\end{tcolorbox}

\begin{tcolorbox}
    \subsection*{Seher \hfill \small{10 \trans{reputation_label}}}
    \tcblower




\begin{tabular}{@{}ll@{}}
    \textbf{Zauberpunkte} & +5\\
    \textbf{Max Arkana} & +3\\
    \textbf{Gewissenhaftigkeit} & +1\\
    \textbf{Auffassungsgabe} & +1\\
    \textbf{Wahrnehmung} & +1\\
    \textbf{Magiekunde} & +1\\
    \textbf{Empathie} & +1\\
    \textbf{Astrale Magie} & True\\
\end{tabular}

\wbif{phasesix}{
    \vspace{3mm}
\faIcon{magic}\hspace{0.5em}}{}

\end{tcolorbox}

\begin{tcolorbox}
    \subsection*{Vertriebsleiter \hfill \small{10 \trans{reputation_label}}}
    \tcblower




\begin{tabular}{@{}ll@{}}
    \textbf{Maximaler Stress} & +1\\
    \textbf{Charme} & +1\\
    \textbf{Überzeugen} & +2\\
    \textbf{Kommunikation} & +2\\
    \textbf{Täuschen und Betrügen} & +1\\
    \textbf{Verwaltung} & +1\\
\end{tabular}

\wbif{phasesix}{
    \vspace{3mm}
\faIcon{plane}\hspace{0.5em}\faIcon{microchip}\hspace{0.5em}}{}

\end{tcolorbox}

\begin{tcolorbox}
    \subsection*{Krankenpfleger \hfill \small{7 \trans{reputation_label}}}
    \tcblower




\begin{tabular}{@{}ll@{}}
    \textbf{Erste Hilfe} & +2\\
    \textbf{Kommunikation} & +1\\
    \textbf{Medizin} & +2\\
\end{tabular}

\wbif{phasesix}{
    \vspace{3mm}
\faIcon{umbrella}\hspace{0.5em}\faIcon{flag}\hspace{0.5em}\faIcon{plane}\hspace{0.5em}\faIcon{microchip}\hspace{0.5em}}{}

\end{tcolorbox}

\begin{tcolorbox}
    \subsection*{Priester \hfill \small{10 \trans{reputation_label}}}
    \tcblower
\textit{Religiös}: Du bist religiös, glaubst an deine Gottheit und verteidigst deinen Glauben auch aktiv. Du darfst Anrufungen aller Art ausführen.


\vspace{3mm}

\begin{tabular}{@{}ll@{}}
    \textbf{Schicksalswürfel} & +1\\
    \textbf{Bildung} & +1\\
    \textbf{Religion} & +2\\
    \textbf{Kommunikation} & +2\\
    \textbf{Musik} & +2\\
\end{tabular}

\wbif{phasesix}{
    \vspace{3mm}
\faIcon{ankh}\hspace{0.5em}}{}

\end{tcolorbox}

\begin{tcolorbox}
    \subsection*{Mönch \hfill \small{7 \trans{reputation_label}}}
    \tcblower
\textit{Religiös}: Du bist religiös, glaubst an deine Gottheit und verteidigst deinen Glauben auch aktiv. Du darfst Anrufungen aller Art ausführen.


\vspace{3mm}

\begin{tabular}{@{}ll@{}}
    \textbf{Religion} & +2\\
    \textbf{Kommunikation} & +1\\
    \textbf{Musik} & +2\\
\end{tabular}

\wbif{phasesix}{
    \vspace{3mm}
\faIcon{ankh}\hspace{0.5em}}{}

\end{tcolorbox}

\begin{tcolorbox}
    \subsection*{Paladin \hfill \small{12 \trans{reputation_label}}}
    \tcblower
\textit{Religiös}: Du bist religiös, glaubst an deine Gottheit und verteidigst deinen Glauben auch aktiv. Du darfst Anrufungen aller Art ausführen.


\vspace{3mm}

\begin{tabular}{@{}ll@{}}
    \textbf{Ausdauer} & +2\\
    \textbf{Resistenz} & +2\\
    \textbf{Nahkampf} & +2\\
    \textbf{Religion} & +3\\
    \textbf{Heraldik} & +1\\
\end{tabular}

\wbif{phasesix}{
    \vspace{3mm}
\faIcon{ankh}\hspace{0.5em}}{}

\end{tcolorbox}

\begin{tcolorbox}
    \subsection*{Vampirjäger \hfill \small{10 \trans{reputation_label}}}
    \tcblower




\begin{tabular}{@{}ll@{}}
    \textbf{Max Arkana} & +3\\
    \textbf{Zauberpunkte} & +5\\
    \textbf{Geschick} & +1\\
    \textbf{Resistenz} & +1\\
    \textbf{Tapferkeit} & +1\\
    \textbf{Heimlichkeit} & +1\\
    \textbf{Zaubern} & +1\\
    \textbf{Weisse Magie} & True\\
\end{tabular}

\wbif{phasesix}{
    \vspace{3mm}
\faIcon{magic}\hspace{0.5em}}{}

\end{tcolorbox}

\begin{tcolorbox}
    \subsection*{Weißmagier \hfill \small{10 \trans{reputation_label}}}
    \tcblower




\begin{tabular}{@{}ll@{}}
    \textbf{Zauberpunkte} & +5\\
    \textbf{Max Arkana} & +3\\
    \textbf{Bildung} & +1\\
    \textbf{Resistenz} & +1\\
    \textbf{Magiekunde} & +1\\
    \textbf{Empathie} & +1\\
    \textbf{Zaubern} & +1\\
    \textbf{Weisse Magie} & True\\
\end{tabular}

\wbif{phasesix}{
    \vspace{3mm}
\faIcon{magic}\hspace{0.5em}}{}

\end{tcolorbox}

\begin{tcolorbox}
    \subsection*{Schwarzmagier \hfill \small{10 \trans{reputation_label}}}
    \tcblower




\begin{tabular}{@{}ll@{}}
    \textbf{Zauberpunkte} & +5\\
    \textbf{Max Arkana} & +3\\
    \textbf{Bildung} & +1\\
    \textbf{Charme} & +1\\
    \textbf{Einschüchtern} & +1\\
    \textbf{Magiekunde} & +1\\
    \textbf{Zaubern} & +1\\
    \textbf{Schwarze Magie} & True\\
\end{tabular}

\wbif{phasesix}{
    \vspace{3mm}
\faIcon{magic}\hspace{0.5em}}{}

\end{tcolorbox}

\begin{tcolorbox}
    \subsection*{Schamane \hfill \small{10 \trans{reputation_label}}}
    \tcblower




\begin{tabular}{@{}ll@{}}
    \textbf{Zauberpunkte} & +5\\
    \textbf{Max Arkana} & +3\\
    \textbf{Willenskraft} & +1\\
    \textbf{Charme} & +1\\
    \textbf{Magiekunde} & +1\\
    \textbf{Empathie} & +1\\
    \textbf{Zaubern} & +1\\
    \textbf{Schamanismus} & True\\
\end{tabular}

\wbif{phasesix}{
    \vspace{3mm}
\faIcon{magic}\hspace{0.5em}}{}

\end{tcolorbox}

\begin{tcolorbox}
    \subsection*{Mystiker \hfill \small{10 \trans{reputation_label}}}
    \tcblower




\begin{tabular}{@{}ll@{}}
    \textbf{Zauberpunkte} & +5\\
    \textbf{Max Arkana} & +3\\
    \textbf{Magiekunde} & +1\\
    \textbf{Darbietung} & +1\\
    \textbf{Zaubern} & +1\\
    \textbf{Mythen und Legenden} & +2\\
    \textbf{Mystizismus} & True\\
\end{tabular}

\wbif{phasesix}{
    \vspace{3mm}
\faIcon{magic}\hspace{0.5em}}{}

\end{tcolorbox}

\begin{tcolorbox}
    \subsection*{Mumienbefrieder \hfill \small{10 \trans{reputation_label}}}
    \tcblower




\begin{tabular}{@{}ll@{}}
    \textbf{Max Arkana} & +3\\
    \textbf{Zauberpunkte} & +5\\
    \textbf{Resistenz} & +1\\
    \textbf{Auffassungsgabe} & +1\\
    \textbf{Zaubern} & +1\\
    \textbf{Magiekunde} & +1\\
    \textbf{Altertümer} & +1\\
    \textbf{Weisse Magie} & True\\
\end{tabular}

\wbif{phasesix}{
    \vspace{3mm}
\faIcon{magic}\hspace{0.5em}}{}

\end{tcolorbox}

\begin{tcolorbox}
    \subsection*{Hermetiker \hfill \small{10 \trans{reputation_label}}}
    \tcblower




\begin{tabular}{@{}ll@{}}
    \textbf{Zauberpunkte} & +5\\
    \textbf{Max Arkana} & +3\\
    \textbf{Bildung} & +1\\
    \textbf{Logik} & +1\\
    \textbf{Überzeugen} & +1\\
    \textbf{Magiekunde} & +1\\
    \textbf{Zaubern} & +1\\
    \textbf{Hermetik} & True\\
\end{tabular}

\wbif{phasesix}{
    \vspace{3mm}
\faIcon{magic}\hspace{0.5em}}{}

\end{tcolorbox}

\begin{tcolorbox}
    \subsection*{Geisterbeschwörer \hfill \small{10 \trans{reputation_label}}}
    \tcblower




\begin{tabular}{@{}ll@{}}
    \textbf{Max Arkana} & +3\\
    \textbf{Zauberpunkte} & +5\\
    \textbf{Charme} & +1\\
    \textbf{Magiekunde} & +1\\
    \textbf{Empathie} & +1\\
    \textbf{Zaubern} & +1\\
    \textbf{Mythen und Legenden} & +1\\
    \textbf{Geisterbeschwörungen} & True\\
\end{tabular}

\wbif{phasesix}{
    \vspace{3mm}
\faIcon{magic}\hspace{0.5em}}{}

\end{tcolorbox}

\begin{tcolorbox}
    \subsection*{Chimärologe \hfill \small{10 \trans{reputation_label}}}
    \tcblower




\begin{tabular}{@{}ll@{}}
    \textbf{Max Arkana} & +3\\
    \textbf{Zauberpunkte} & +5\\
    \textbf{Resistenz} & +1\\
    \textbf{Kraft} & +1\\
    \textbf{Tapferkeit} & +1\\
    \textbf{Magiekunde} & +1\\
    \textbf{Zaubern} & +1\\
    \textbf{Chimärologie} & True\\
\end{tabular}

\wbif{phasesix}{
    \vspace{3mm}
\faIcon{magic}\hspace{0.5em}}{}

\end{tcolorbox}

\begin{tcolorbox}
    \subsection*{Astrologe \hfill \small{10 \trans{reputation_label}}}
    \tcblower




\begin{tabular}{@{}ll@{}}
    \textbf{Zauberpunkte} & +5\\
    \textbf{Max Arkana} & +3\\
    \textbf{Auffassungsgabe} & +1\\
    \textbf{Gewissenhaftigkeit} & +1\\
    \textbf{Magiekunde} & +1\\
    \textbf{Zaubern} & +1\\
    \textbf{Mechanik} & +1\\
    \textbf{Astrale Magie} & True\\
\end{tabular}

\wbif{phasesix}{
    \vspace{3mm}
\faIcon{magic}\hspace{0.5em}}{}

\end{tcolorbox}

\begin{tcolorbox}
    \subsection*{Medium \hfill \small{8 \trans{reputation_label}}}
    \tcblower




\begin{tabular}{@{}ll@{}}
    \textbf{Maximaler Stress} & +1\\
    \textbf{Empathie} & +2\\
    \textbf{Mythen und Legenden} & +2\\
\end{tabular}

\wbif{phasesix}{
    \vspace{3mm}
\faIcon{cube}\hspace{0.5em}}{}

\end{tcolorbox}

\begin{tcolorbox}
    \subsection*{Arzt \hfill \small{11 \trans{reputation_label}}}
    \tcblower
Ein Arzt ist ein Spezialist wenn es um die Gesundheit des Körpers und/oder Geistes des Menschen geht.

\begin{pquote}{None}
"We need a Doctor!"
\end{pquote}

\vspace{3mm}

\begin{tabular}{@{}ll@{}}
    \textbf{Erste Hilfe} & +3\\
    \textbf{Untersuchen} & +1\\
    \textbf{Medizin} & +4\\
    \textbf{Allgemeinwissen} & +1\\
\end{tabular}

\wbif{phasesix}{
    \vspace{3mm}
\faIcon{umbrella}\hspace{0.5em}\faIcon{flag}\hspace{0.5em}\faIcon{plane}\hspace{0.5em}\faIcon{microchip}\hspace{0.5em}}{}

\end{tcolorbox}

\begin{tcolorbox}
    \subsection*{Jäger \hfill \small{9 \trans{reputation_label}}}
    \tcblower




\begin{tabular}{@{}ll@{}}
    \textbf{Geschick} & +1\\
    \textbf{Schnelligkeit} & +2\\
    \textbf{Schießen} & +2\\
    \textbf{Wahrnehmung} & +2\\
\end{tabular}

\wbif{phasesix}{
    \vspace{3mm}
\faIcon{chess-rook}\hspace{0.5em}\faIcon{dragon}\hspace{0.5em}\faIcon{ankh}\hspace{0.5em}\faIcon{magic}\hspace{0.5em}}{}

\end{tcolorbox}

\section{Bildung}

Die Charakterschablonen der Bildung stellen die Schulbildung des Charakters dar.

\begin{tcolorbox}
    \subsection*{Hauptschule \hfill \small{4 \trans{reputation_label}}}
    \tcblower




\begin{tabular}{@{}ll@{}}
    \textbf{Schicksalswürfel} & +1\\
    \textbf{Resistenz} & +1\\
    \textbf{Einschüchtern} & +1\\
\end{tabular}

\wbif{phasesix}{
    \vspace{3mm}
\faIcon{flag}\hspace{0.5em}\faIcon{plane}\hspace{0.5em}\faIcon{microchip}\hspace{0.5em}}{}

\end{tcolorbox}

\begin{tcolorbox}
    \subsection*{Abitur \hfill \small{6 \trans{reputation_label}}}
    \tcblower




\begin{tabular}{@{}ll@{}}
    \textbf{Wiederholungswürfe} & +1\\
    \textbf{Bildung} & +2\\
    \textbf{Logik} & +1\\
\end{tabular}

\wbif{phasesix}{
    \vspace{3mm}
\faIcon{flag}\hspace{0.5em}\faIcon{plane}\hspace{0.5em}\faIcon{microchip}\hspace{0.5em}}{}

\end{tcolorbox}

\begin{tcolorbox}
    \subsection*{Schule abgebrochen \hfill \small{-1 \trans{reputation_label}}}
    \tcblower


\begin{pquote}{Jacky Tyler}
There's no point in gettin' up sweetheart. There is no job to go to.
\end{pquote}

\vspace{3mm}

\begin{tabular}{@{}ll@{}}
    \textbf{Gewissenhaftigkeit} & -2\\
    \textbf{Nahkampf} & +1\\
\end{tabular}

\wbif{phasesix}{
    \vspace{3mm}
\faIcon{umbrella}\hspace{0.5em}\faIcon{flag}\hspace{0.5em}\faIcon{plane}\hspace{0.5em}\faIcon{microchip}\hspace{0.5em}\faIcon{space-shuttle}\hspace{0.5em}}{}

\end{tcolorbox}

\begin{tcolorbox}
    \subsection*{Gesamtschule \hfill \small{6 \trans{reputation_label}}}
    \tcblower




\begin{tabular}{@{}ll@{}}
    \textbf{Bildung} & +1\\
    \textbf{Mechanik} & +2\\
    \textbf{Kommunikation} & +1\\
\end{tabular}

\wbif{phasesix}{
    \vspace{3mm}
\faIcon{plane}\hspace{0.5em}\faIcon{microchip}\hspace{0.5em}}{}

\end{tcolorbox}

\begin{tcolorbox}
    \subsection*{Internat \hfill \small{10 \trans{reputation_label}}}
    \tcblower
\textit{Vermögen}: Deine Familie hat ein nennenswertes Vermögen angehäuft, von dem du noch viele Jahre bequem zehren kannst.


\vspace{3mm}

\begin{tabular}{@{}ll@{}}
    \textbf{Auffassungsgabe} & +2\\
    \textbf{Bildung} & +2\\
    \textbf{Kommunikation} & +2\\
\end{tabular}

\wbif{phasesix}{
    \vspace{3mm}
\faIcon{umbrella}\hspace{0.5em}\faIcon{flag}\hspace{0.5em}\faIcon{plane}\hspace{0.5em}\faIcon{microchip}\hspace{0.5em}}{}

\end{tcolorbox}

\begin{tcolorbox}
    \subsection*{Youtube \hfill \small{3 \trans{reputation_label}}}
    \tcblower


\begin{pquote}{Axel Stoll}
Physik ist Magie durch Wollen
\end{pquote}

\vspace{3mm}

\begin{tabular}{@{}ll@{}}
    \textbf{Bildung} & -1\\
    \textbf{Untersuchen} & +1\\
    \textbf{Kommunikation} & +1\\
\end{tabular}

\wbif{phasesix}{
    \vspace{3mm}
\faIcon{microchip}\hspace{0.5em}}{}

\end{tcolorbox}

\begin{tcolorbox}
    \subsection*{Hausunterricht \hfill \small{0 \trans{reputation_label}}}
    \tcblower
\textit{Religiös}: Du bist religiös, glaubst an deine Gottheit und verteidigst deinen Glauben auch aktiv.


\vspace{3mm}

\begin{tabular}{@{}ll@{}}
    \textbf{Bildung} & -2\\
    \textbf{Auffassungsgabe} & -1\\
    \textbf{Religion} & +2\\
\end{tabular}

\wbif{phasesix}{
    \vspace{3mm}
\faIcon{eye}\hspace{0.5em}\faIcon{umbrella}\hspace{0.5em}\faIcon{flag}\hspace{0.5em}\faIcon{plane}\hspace{0.5em}\faIcon{microchip}\hspace{0.5em}}{}

\end{tcolorbox}

\begin{tcolorbox}
    \subsection*{Freisprechung \hfill \small{10 \trans{reputation_label}}}
    \tcblower




\begin{tabular}{@{}ll@{}}
    \textbf{Bonuswürfel} & +2\\
    \textbf{Geschick} & +1\\
    \textbf{Akrobatik} & +1\\
    \textbf{Fahren} & +1\\
    \textbf{Mechanik} & +2\\
\end{tabular}

\wbif{phasesix}{
    \vspace{3mm}
\faIcon{chess-rook}\hspace{0.5em}}{}

\end{tcolorbox}

\begin{tcolorbox}
    \subsection*{Trivium \hfill \small{12 \trans{reputation_label}}}
    \tcblower




\begin{tabular}{@{}ll@{}}
    \textbf{Bonuswürfel} & +1\\
    \textbf{Bildung} & +2\\
    \textbf{Kommunikation} & +3\\
    \textbf{Altertümer} & +1\\
    \textbf{Lesen/Schreiben} & +2\\
    \textbf{Rhetorik} & +2\\
\end{tabular}

\wbif{phasesix}{
    \vspace{3mm}
\faIcon{chess-rook}\hspace{0.5em}}{}

\end{tcolorbox}

\begin{tcolorbox}
    \subsection*{Quadrivium \hfill \small{18 \trans{reputation_label}}}
    \tcblower




\begin{tabular}{@{}ll@{}}
    \textbf{Schicksalswürfel} & +2\\
    \textbf{Bildung} & +2\\
    \textbf{Historie} & +2\\
    \textbf{Politik} & +1\\
    \textbf{Astronomie} & +2\\
    \textbf{Musik} & +2\\
    \textbf{Mathematik} & +2\\
    \textbf{Lesen/Schreiben} & +2\\
\end{tabular}

\wbif{phasesix}{
    \vspace{3mm}
\faIcon{chess-rook}\hspace{0.5em}}{}

\end{tcolorbox}

\begin{tcolorbox}
    \subsection*{Studium \hfill \small{10 \trans{reputation_label}}}
    \tcblower




\begin{tabular}{@{}ll@{}}
    \textbf{Bonuswürfel} & +1\\
    \textbf{Bildung} & +2\\
    \textbf{Gewissenhaftigkeit} & +1\\
    \textbf{Logik} & +1\\
    \textbf{Naturkunde} & +1\\
    \textbf{Historie} & +1\\
    \textbf{Kommunikation} & +2\\
\end{tabular}

\wbif{phasesix}{
    \vspace{3mm}
\faIcon{eye}\hspace{0.5em}\faIcon{umbrella}\hspace{0.5em}\faIcon{flag}\hspace{0.5em}\faIcon{plane}\hspace{0.5em}\faIcon{microchip}\hspace{0.5em}}{}

\end{tcolorbox}

\begin{tcolorbox}
    \subsection*{Militärakademie \hfill \small{11 \trans{reputation_label}}}
    \tcblower


\begin{pquote}{Douglas MacArthur}
Whoever said the pen is mightier than the sword obviously never encountered automatic weapons.
\end{pquote}

\vspace{3mm}

\begin{tabular}{@{}ll@{}}
    \textbf{Schicksalswürfel} & +1\\
    \textbf{Nahkampf} & +1\\
    \textbf{Einschüchtern} & +1\\
    \textbf{Erste Hilfe} & +1\\
    \textbf{Politik} & +1\\
    \textbf{Schießen} & +1\\
    \textbf{Kriegsführung} & +2\\
    \textbf{Lesen/Schreiben} & +1\\
    \textbf{Reiten} & +1\\
\end{tabular}

\wbif{phasesix}{
    \vspace{3mm}
\faIcon{cube}\hspace{0.5em}}{}

\end{tcolorbox}

\begin{tcolorbox}
    \subsection*{Arkanes Studium \hfill \small{10 \trans{reputation_label}}}
    \tcblower




\begin{tabular}{@{}ll@{}}
    \textbf{Max Arkana} & +1\\
    \textbf{Zauberpunkte} & +5\\
    \textbf{Logik} & +1\\
    \textbf{Willenskraft} & +1\\
    \textbf{Zaubern} & +1\\
    \textbf{Magiekunde} & +2\\
\end{tabular}

\wbif{phasesix}{
    \vspace{3mm}
\faIcon{magic}\hspace{0.5em}}{}

\end{tcolorbox}

\begin{tcolorbox}
    \subsection*{Arkaner Lehrmeister \hfill \small{10 \trans{reputation_label}}}
    \tcblower


\begin{pquote}{Otfried Preußler}
Es gibt eine Art von Zauberei, die man mühsam erlernen muß: Das ist die, wie sie im Koraktor steht, Zeichen für Zeichen und Formel um Formel. Und dann gibt es eine, die wächst einem aus der Tiefe des Herzens zu: aus der Sorge um jemanden, den man lieb hat. Ich weiß, daß das schwer zu begreifen ist - aber du solltest darauf vertrauen, Krabat.
\end{pquote}

\vspace{3mm}

\begin{tabular}{@{}ll@{}}
    \textbf{Max Arkana} & +1\\
    \textbf{Zauberpunkte} & +15\\
    \textbf{Willenskraft} & +1\\
    \textbf{Zaubern} & +2\\
    \textbf{Magiekunde} & +1\\
\end{tabular}

\wbif{phasesix}{
    \vspace{3mm}
\faIcon{magic}\hspace{0.5em}}{}

\end{tcolorbox}

\begin{tcolorbox}
    \subsection*{Arkane Schule \hfill \small{10 \trans{reputation_label}}}
    \tcblower




\begin{tabular}{@{}ll@{}}
    \textbf{Zauberpunkte} & +10\\
    \textbf{Max Arkana} & +2\\
    \textbf{Bildung} & +1\\
    \textbf{Zaubern} & +1\\
    \textbf{Magiekunde} & +2\\
\end{tabular}

\wbif{phasesix}{
    \vspace{3mm}
\faIcon{magic}\hspace{0.5em}}{}

\end{tcolorbox}

\begin{tcolorbox}
    \subsection*{Fachidiot \hfill \small{0 \trans{reputation_label}}}
    \tcblower
In einem engen Fachgebiet deiner Wahl würfelst du Wissensproben mit 3 zusätzlichen Würfeln.

\begin{pquote}{Oscar Wilde}
Der Experte ist ein gewöhnlicher Mann, der - wenn er nicht daheim ist - Ratschläge erteilt.
\end{pquote}

\vspace{3mm}

\begin{tabular}{@{}ll@{}}
    \textbf{Bildung} & -1\\
    \textbf{Gewissenhaftigkeit} & -1\\
\end{tabular}

\wbif{phasesix}{
    \vspace{3mm}
\faIcon{plane}\hspace{0.5em}\faIcon{microchip}\hspace{0.5em}}{}

\end{tcolorbox}

\begin{tcolorbox}
    \subsection*{Arkane Meditation \hfill \small{6 \trans{reputation_label}}}
    \tcblower




\begin{tabular}{@{}ll@{}}
    \textbf{Zauberpunkte} & +5\\
    \textbf{Max Arkana} & +2\\
    \textbf{Magiekunde} & +2\\
\end{tabular}

\wbif{phasesix}{
    \vspace{3mm}
\faIcon{magic}\hspace{0.5em}}{}

\end{tcolorbox}

\begin{tcolorbox}
    \subsection*{Schildknappe \hfill \small{5 \trans{reputation_label}}}
    \tcblower


\begin{pquote}{Miguel de Cervantes Saavedra}
Wer ein guter Schildknappe gewesen ist, wird auch ein guter Statthalter sein.
\end{pquote}

\vspace{3mm}

\begin{tabular}{@{}ll@{}}
    \textbf{Ausdauer} & +1\\
    \textbf{Resistenz} & +1\\
    \textbf{Nahkampf} & +1\\
\end{tabular}

\wbif{phasesix}{
    \vspace{3mm}
\faIcon{chess-rook}\hspace{0.5em}}{}

\end{tcolorbox}

\section{Interessen}

Diese Charakterschablonen beschreiben die Interessen des Charakters. Sie verändern ausgewählte Eigenschaften, sind dabei aber günstiger und verändern weniger als Berufe.

\begin{tcolorbox}
    \subsection*{Sport \hfill \small{7 \trans{reputation_label}}}
    \tcblower




\begin{tabular}{@{}ll@{}}
    \textbf{Schnelligkeit} & +1\\
    \textbf{Ausdauer} & +2\\
    \textbf{Werfen} & +1\\
    \textbf{Akrobatik} & +1\\
\end{tabular}

\wbif{phasesix}{
    \vspace{3mm}
\faIcon{flag}\hspace{0.5em}\faIcon{plane}\hspace{0.5em}\faIcon{microchip}\hspace{0.5em}}{}

\end{tcolorbox}

\begin{tcolorbox}
    \subsection*{Jagd \hfill \small{6 \trans{reputation_label}}}
    \tcblower




\begin{tabular}{@{}ll@{}}
    \textbf{Ausdauer} & +1\\
    \textbf{Schießen} & +2\\
    \textbf{Zoologie} & +1\\
\end{tabular}

\wbif{phasesix}{
    \vspace{3mm}
\faIcon{cube}\hspace{0.5em}}{}

\end{tcolorbox}

\begin{tcolorbox}
    \subsection*{Musik \hfill \small{6 \trans{reputation_label}}}
    \tcblower




\begin{tabular}{@{}ll@{}}
    \textbf{Geschick} & +1\\
    \textbf{Darbietung} & +1\\
    \textbf{Musik} & +2\\
\end{tabular}

\wbif{phasesix}{
    \vspace{3mm}
\faIcon{cube}\hspace{0.5em}}{}

\end{tcolorbox}

\begin{tcolorbox}
    \subsection*{Lesen \hfill \small{3 \trans{reputation_label}}}
    \tcblower




\begin{tabular}{@{}ll@{}}
    \textbf{Bildung} & +1\\
    \textbf{Allgemeinwissen} & +2\\
\end{tabular}

\wbif{phasesix}{
    \vspace{3mm}
\faIcon{umbrella}\hspace{0.5em}\faIcon{flag}\hspace{0.5em}\faIcon{plane}\hspace{0.5em}\faIcon{microchip}\hspace{0.5em}\faIcon{space-shuttle}\hspace{0.5em}}{}

\end{tcolorbox}

\begin{tcolorbox}
    \subsection*{Kunst \hfill \small{2 \trans{reputation_label}}}
    \tcblower




\begin{tabular}{@{}ll@{}}
    \textbf{Darbietung} & +2\\
\end{tabular}

\wbif{phasesix}{
    \vspace{3mm}
\faIcon{cube}\hspace{0.5em}}{}

\end{tcolorbox}

\begin{tcolorbox}
    \subsection*{Trekking \hfill \small{5 \trans{reputation_label}}}
    \tcblower




\begin{tabular}{@{}ll@{}}
    \textbf{Ausdauer} & +2\\
    \textbf{Erste Hilfe} & +1\\
    \textbf{Orientierung} & +2\\
\end{tabular}

\wbif{phasesix}{
    \vspace{3mm}
\faIcon{plane}\hspace{0.5em}\faIcon{microchip}\hspace{0.5em}}{}

\end{tcolorbox}

\begin{tcolorbox}
    \subsection*{Karate \hfill \small{3 \trans{reputation_label}}}
    \tcblower




\begin{tabular}{@{}ll@{}}
    \textbf{Geschick} & +1\\
    \textbf{Nahkampf} & +2\\
\end{tabular}

\wbif{phasesix}{
    \vspace{3mm}
\faIcon{plane}\hspace{0.5em}\faIcon{microchip}\hspace{0.5em}}{}

\end{tcolorbox}

\begin{tcolorbox}
    \subsection*{Lesen \hfill \small{5 \trans{reputation_label}}}
    \tcblower




\begin{tabular}{@{}ll@{}}
    \textbf{Gewissenhaftigkeit} & +1\\
    \textbf{Lesen/Schreiben} & +2\\
\end{tabular}

\wbif{phasesix}{
    \vspace{3mm}
\faIcon{eye}\hspace{0.5em}\faIcon{chess-rook}\hspace{0.5em}}{}

\end{tcolorbox}

\begin{tcolorbox}
    \subsection*{Geschichte \hfill \small{4 \trans{reputation_label}}}
    \tcblower




\begin{tabular}{@{}ll@{}}
    \textbf{Historie} & +2\\
    \textbf{Altertümer} & +2\\
\end{tabular}

\wbif{phasesix}{
    \vspace{3mm}
\faIcon{cube}\hspace{0.5em}}{}

\end{tcolorbox}

\begin{tcolorbox}
    \subsection*{Handarbeit \hfill \small{4 \trans{reputation_label}}}
    \tcblower




\begin{tabular}{@{}ll@{}}
    \textbf{Geschick} & +2\\
\end{tabular}

\wbif{phasesix}{
    \vspace{3mm}
\faIcon{cube}\hspace{0.5em}}{}

\end{tcolorbox}

\begin{tcolorbox}
    \subsection*{Esoterik \hfill \small{5 \trans{reputation_label}}}
    \tcblower




\begin{tabular}{@{}ll@{}}
    \textbf{Schicksalswürfel} & +1\\
    \textbf{Logik} & -2\\
    \textbf{Einschüchtern} & +1\\
    \textbf{Heimlichkeit} & +1\\
    \textbf{Täuschen und Betrügen} & +1\\
\end{tabular}

\wbif{phasesix}{
    \vspace{3mm}
\faIcon{cube}\hspace{0.5em}}{}

\end{tcolorbox}

\begin{tcolorbox}
    \subsection*{Schützenverein \hfill \small{3 \trans{reputation_label}}}
    \tcblower




\begin{tabular}{@{}ll@{}}
    \textbf{Bonuswürfel} & +1\\
    \textbf{Attraktivität} & -1\\
    \textbf{Schießen} & +1\\
\end{tabular}

\wbif{phasesix}{
    \vspace{3mm}
\faIcon{flag}\hspace{0.5em}\faIcon{plane}\hspace{0.5em}\faIcon{microchip}\hspace{0.5em}}{}

\end{tcolorbox}

\begin{tcolorbox}
    \subsection*{P\&P Rollenspiele \hfill \small{6 \trans{reputation_label}}}
    \tcblower




\begin{tabular}{@{}ll@{}}
    \textbf{Historie} & +1\\
    \textbf{Kommunikation} & +2\\
    \textbf{Darbietung} & +1\\
\end{tabular}

\wbif{phasesix}{
    \vspace{3mm}
\faIcon{plane}\hspace{0.5em}\faIcon{microchip}\hspace{0.5em}}{}

\end{tcolorbox}

\begin{tcolorbox}
    \subsection*{Briefmarken sammeln \hfill \small{4 \trans{reputation_label}}}
    \tcblower




\begin{tabular}{@{}ll@{}}
    \textbf{Gewissenhaftigkeit} & +2\\
\end{tabular}

\wbif{phasesix}{
    \vspace{3mm}
\faIcon{plane}\hspace{0.5em}\faIcon{microchip}\hspace{0.5em}}{}

\end{tcolorbox}

\begin{tcolorbox}
    \subsection*{Heraldik \hfill \small{8 \trans{reputation_label}}}
    \tcblower




\begin{tabular}{@{}ll@{}}
    \textbf{Gewissenhaftigkeit} & +2\\
    \textbf{Wahrnehmung} & +1\\
    \textbf{Historie} & +1\\
    \textbf{Heraldik} & +2\\
\end{tabular}

\wbif{phasesix}{
    \vspace{3mm}
\faIcon{chess-rook}\hspace{0.5em}}{}

\end{tcolorbox}

\begin{tcolorbox}
    \subsection*{Yoga \hfill \small{6 \trans{reputation_label}}}
    \tcblower




\begin{tabular}{@{}ll@{}}
    \textbf{Geschick} & +2\\
    \textbf{Ausdauer} & +1\\
    \textbf{Akrobatik} & +1\\
\end{tabular}

\wbif{phasesix}{
    \vspace{3mm}
\faIcon{microchip}\hspace{0.5em}}{}

\end{tcolorbox}

\begin{tcolorbox}
    \subsection*{Rettungsschwimmen \hfill \small{7 \trans{reputation_label}}}
    \tcblower




\begin{tabular}{@{}ll@{}}
    \textbf{Ausdauer} & +2\\
    \textbf{Erste Hilfe} & +2\\
    \textbf{Tapferkeit} & +1\\
\end{tabular}

\wbif{phasesix}{
    \vspace{3mm}
\faIcon{plane}\hspace{0.5em}\faIcon{microchip}\hspace{0.5em}}{}

\end{tcolorbox}

\begin{tcolorbox}
    \subsection*{Kultmitgliedschaft \hfill \small{5 \trans{reputation_label}}}
    \tcblower




\begin{tabular}{@{}ll@{}}
    \textbf{Wiederholungswürfe} & +1\\
    \textbf{Schicksalswürfel} & +1\\
    \textbf{Bonuswürfel} & -2\\
\end{tabular}

\wbif{phasesix}{
    \vspace{3mm}
\faIcon{cube}\hspace{0.5em}}{}

\end{tcolorbox}

\begin{tcolorbox}
    \subsection*{Burschenschaft \hfill \small{4 \trans{reputation_label}}}
    \tcblower




\begin{tabular}{@{}ll@{}}
    \textbf{Charme} & +1\\
    \textbf{Attraktivität} & -1\\
    \textbf{Nahkampf} & +1\\
    \textbf{Historie} & +1\\
    \textbf{Etikette} & +2\\
\end{tabular}

\wbif{phasesix}{
    \vspace{3mm}
\faIcon{umbrella}\hspace{0.5em}\faIcon{flag}\hspace{0.5em}\faIcon{plane}\hspace{0.5em}\faIcon{microchip}\hspace{0.5em}}{}

\end{tcolorbox}

\begin{tcolorbox}
    \subsection*{Tanzen \hfill \small{6 \trans{reputation_label}}}
    \tcblower




\begin{tabular}{@{}ll@{}}
    \textbf{Geschick} & +1\\
    \textbf{Ausdauer} & +1\\
    \textbf{Attraktivität} & +2\\
\end{tabular}

\wbif{phasesix}{
    \vspace{3mm}
\faIcon{cube}\hspace{0.5em}}{}

\end{tcolorbox}

\begin{tcolorbox}
    \subsection*{Fahrzeugtuning \hfill \small{5 \trans{reputation_label}}}
    \tcblower




\begin{tabular}{@{}ll@{}}
    \textbf{Fahren} & +1\\
    \textbf{Mechanik} & +1\\
    \textbf{Fahrzeuge} & +1\\
\end{tabular}

\wbif{phasesix}{
    \vspace{3mm}
\faIcon{plane}\hspace{0.5em}\faIcon{microchip}\hspace{0.5em}\faIcon{space-shuttle}\hspace{0.5em}}{}

\end{tcolorbox}

\begin{tcolorbox}
    \subsection*{Chemie \hfill \small{6 \trans{reputation_label}}}
    \tcblower




\begin{tabular}{@{}ll@{}}
    \textbf{Gewissenhaftigkeit} & +1\\
    \textbf{Tapferkeit} & +1\\
    \textbf{Chemie} & +2\\
\end{tabular}

\wbif{phasesix}{
    \vspace{3mm}
\faIcon{umbrella}\hspace{0.5em}\faIcon{flag}\hspace{0.5em}\faIcon{plane}\hspace{0.5em}\faIcon{microchip}\hspace{0.5em}}{}

\end{tcolorbox}

\begin{tcolorbox}
    \subsection*{Reiten \hfill \small{5 \trans{reputation_label}}}
    \tcblower




\begin{tabular}{@{}ll@{}}
    \textbf{Fahren} & +1\\
    \textbf{Reiten} & +2\\
\end{tabular}

\wbif{phasesix}{
    \vspace{3mm}
\faIcon{cube}\hspace{0.5em}}{}

\end{tcolorbox}

\begin{tcolorbox}
    \subsection*{Graffiti sprayen \hfill \small{4 \trans{reputation_label}}}
    \tcblower




\begin{tabular}{@{}ll@{}}
    \textbf{Heimlichkeit} & +1\\
    \textbf{Überzeugen} & +1\\
\end{tabular}

\wbif{phasesix}{
    \vspace{3mm}
\faIcon{plane}\hspace{0.5em}\faIcon{microchip}\hspace{0.5em}}{}

\end{tcolorbox}

\begin{tcolorbox}
    \subsection*{Parkour \hfill \small{7 \trans{reputation_label}}}
    \tcblower




\begin{tabular}{@{}ll@{}}
    \textbf{Schnelligkeit} & +1\\
    \textbf{Geschick} & +1\\
    \textbf{Ausdauer} & +1\\
    \textbf{Akrobatik} & +2\\
\end{tabular}

\wbif{phasesix}{
    \vspace{3mm}
\faIcon{microchip}\hspace{0.5em}}{}

\end{tcolorbox}

\begin{tcolorbox}
    \subsection*{Workaholism \hfill \small{-4 \trans{reputation_label}}}
    \tcblower




\begin{tabular}{@{}ll@{}}
    \textbf{Logik} & -1\\
\end{tabular}

\wbif{phasesix}{
    \vspace{3mm}
\faIcon{microchip}\hspace{0.5em}}{}

\end{tcolorbox}

\begin{tcolorbox}
    \subsection*{Alchemie \hfill \small{5 \trans{reputation_label}}}
    \tcblower


\begin{pquote}{Paracelsus}
Alles ist Gift, ausschlaggebend ist nur die Menge. Alles Tun ist ein alchemistisches Zuendeführen, eine geistige Goldmachung und Kunst der Vollendung. Alles Wachsen ist Auferstehen. Auch in die Liebe muss man hineinwachsen und ihre Stunden abwarten, denn die Gewächse der Erde und die Gaben im Menschen haben ihre Zeit.
\end{pquote}

\vspace{3mm}

\begin{tabular}{@{}ll@{}}
    \textbf{Naturkunde} & +2\\
    \textbf{Alchemie} & +1\\
\end{tabular}

\wbif{phasesix}{
    \vspace{3mm}
\faIcon{chess-rook}\hspace{0.5em}}{}

\end{tcolorbox}

\begin{tcolorbox}
    \subsection*{Sudoku \hfill \small{3 \trans{reputation_label}}}
    \tcblower




\begin{tabular}{@{}ll@{}}
    \textbf{Logik} & +1\\
\end{tabular}

\wbif{phasesix}{
    \vspace{3mm}
\faIcon{plane}\hspace{0.5em}\faIcon{microchip}\hspace{0.5em}}{}

\end{tcolorbox}

\begin{tcolorbox}
    \subsection*{Kochen \hfill \small{3 \trans{reputation_label}}}
    \tcblower


\begin{pquote}{Gordon Ramsay}
My gran could do better! And she’s dead!
\end{pquote}

\vspace{3mm}

\begin{tabular}{@{}ll@{}}
    \textbf{Kochen} & +2\\
\end{tabular}

\wbif{phasesix}{
    \vspace{3mm}
\faIcon{cube}\hspace{0.5em}}{}

\end{tcolorbox}

\begin{tcolorbox}
    \subsection*{Arkanes Training \hfill \small{3 \trans{reputation_label}}}
    \tcblower




\begin{tabular}{@{}ll@{}}
    \textbf{Zauberpunkte} & +5\\
\end{tabular}

\wbif{phasesix}{
    \vspace{3mm}
\faIcon{magic}\hspace{0.5em}}{}

\end{tcolorbox}

\begin{tcolorbox}
    \subsection*{Wissbegierig \hfill \small{3 \trans{reputation_label}}}
    \tcblower




\begin{tabular}{@{}ll@{}}
    \textbf{Bildung} & +2\\
\end{tabular}

\wbif{phasesix}{
    \vspace{3mm}
\faIcon{cube}\hspace{0.5em}}{}

\end{tcolorbox}

\begin{tcolorbox}
    \subsection*{Krafttraining \hfill \small{3 \trans{reputation_label}}}
    \tcblower




\begin{tabular}{@{}ll@{}}
    \textbf{Kraft} & +1\\
\end{tabular}

\wbif{phasesix}{
    \vspace{3mm}
\faIcon{cube}\hspace{0.5em}}{}

\end{tcolorbox}

\begin{tcolorbox}
    \subsection*{Akolyth \hfill \small{5 \trans{reputation_label}}}
    \tcblower


\begin{pquote}{Melissa Marr}
He had hopes, but hope wasn't a solution.
\end{pquote}

\vspace{3mm}

\begin{tabular}{@{}ll@{}}
    \textbf{Religion} & +2\\
    \textbf{Etikette} & +1\\
\end{tabular}

\wbif{phasesix}{
    \vspace{3mm}
\faIcon{ankh}\hspace{0.5em}}{}

\end{tcolorbox}

\section{Charakter}

Diese Charakterschablonen beschreiben bestimmte Charaktereigenschaften. Sie verändern einige wenige Eigenschaften des Charakters, fügen Fertigkeiten und Wissen hinzu und sind im Allgemeinen günstiger als Berufe.

\begin{tcolorbox}
    \subsection*{Tausendsassa \hfill \small{6 \trans{reputation_label}}}
    \tcblower




\begin{tabular}{@{}ll@{}}
    \textbf{Bonuswürfel} & +2\\
    \textbf{Wiederholungswürfe} & +1\\
    \textbf{Kommunikation} & +1\\
\end{tabular}

\wbif{phasesix}{
    \vspace{3mm}
\faIcon{cube}\hspace{0.5em}}{}

\end{tcolorbox}

\begin{tcolorbox}
    \subsection*{Spieler \hfill \small{-5 \trans{reputation_label}}}
    \tcblower




\begin{tabular}{@{}ll@{}}
    \textbf{Gewissenhaftigkeit} & -2\\
\end{tabular}

\wbif{phasesix}{
    \vspace{3mm}
\faIcon{cube}\hspace{0.5em}}{}

\end{tcolorbox}

\begin{tcolorbox}
    \subsection*{Raucher \hfill \small{-5 \trans{reputation_label}}}
    \tcblower




\begin{tabular}{@{}ll@{}}
    \textbf{Ausdauer} & -2\\
    \textbf{Attraktivität} & -1\\
\end{tabular}

\wbif{phasesix}{
    \vspace{3mm}
\faIcon{cube}\hspace{0.5em}}{}

\end{tcolorbox}

\begin{tcolorbox}
    \subsection*{Robust \hfill \small{4 \trans{reputation_label}}}
    \tcblower




\begin{tabular}{@{}ll@{}}
    \textbf{Maximale Lebenspunkte} & +2\\
\end{tabular}

\wbif{phasesix}{
    \vspace{3mm}
\faIcon{cube}\hspace{0.5em}}{}

\end{tcolorbox}

\begin{tcolorbox}
    \subsection*{Notorischer Lügner \hfill \small{2 \trans{reputation_label}}}
    \tcblower




\begin{tabular}{@{}ll@{}}
    \textbf{Logik} & -1\\
    \textbf{Täuschen und Betrügen} & +2\\
\end{tabular}

\wbif{phasesix}{
    \vspace{3mm}
\faIcon{cube}\hspace{0.5em}}{}

\end{tcolorbox}

\begin{tcolorbox}
    \subsection*{Philanthrop \hfill \small{4 \trans{reputation_label}}}
    \tcblower




\begin{tabular}{@{}ll@{}}
    \textbf{Attraktivität} & +1\\
    \textbf{Empathie} & +1\\
\end{tabular}

\wbif{phasesix}{
    \vspace{3mm}
\faIcon{cube}\hspace{0.5em}}{}

\end{tcolorbox}

\begin{tcolorbox}
    \subsection*{Korrupt \hfill \small{3 \trans{reputation_label}}}
    \tcblower




\begin{tabular}{@{}ll@{}}
    \textbf{Gewissenhaftigkeit} & -2\\
    \textbf{Täuschen und Betrügen} & +2\\
    \textbf{Einschüchtern} & +1\\
\end{tabular}

\wbif{phasesix}{
    \vspace{3mm}
\faIcon{cube}\hspace{0.5em}}{}

\end{tcolorbox}

\begin{tcolorbox}
    \subsection*{Säufer \hfill \small{-5 \trans{reputation_label}}}
    \tcblower




\begin{tabular}{@{}ll@{}}
    \textbf{Schicksalswürfel} & +1\\
    \textbf{Auffassungsgabe} & -2\\
    \textbf{Wahrnehmung} & -1\\
\end{tabular}

\wbif{phasesix}{
    \vspace{3mm}
\faIcon{cube}\hspace{0.5em}}{}

\end{tcolorbox}

\begin{tcolorbox}
    \subsection*{Betrüger \hfill \small{5 \trans{reputation_label}}}
    \tcblower




\begin{tabular}{@{}ll@{}}
    \textbf{Täuschen und Betrügen} & +3\\
\end{tabular}

\wbif{phasesix}{
    \vspace{3mm}
\faIcon{cube}\hspace{0.5em}}{}

\end{tcolorbox}

\begin{tcolorbox}
    \subsection*{Chauvinist \hfill \small{2 \trans{reputation_label}}}
    \tcblower




\begin{tabular}{@{}ll@{}}
    \textbf{Charme} & -2\\
    \textbf{Attraktivität} & +2\\
\end{tabular}

\wbif{phasesix}{
    \vspace{3mm}
\faIcon{cube}\hspace{0.5em}}{}

\end{tcolorbox}

\begin{tcolorbox}
    \subsection*{Sympathisch \hfill \small{5 \trans{reputation_label}}}
    \tcblower




\begin{tabular}{@{}ll@{}}
    \textbf{Attraktivität} & +2\\
    \textbf{Charme} & +1\\
\end{tabular}

\wbif{phasesix}{
    \vspace{3mm}
\faIcon{cube}\hspace{0.5em}}{}

\end{tcolorbox}

\begin{tcolorbox}
    \subsection*{Risikofreudig \hfill \small{2 \trans{reputation_label}}}
    \tcblower




\begin{tabular}{@{}ll@{}}
    \textbf{Gewissenhaftigkeit} & -2\\
    \textbf{Resistenz} & +1\\
    \textbf{Auffassungsgabe} & +1\\
\end{tabular}

\wbif{phasesix}{
    \vspace{3mm}
\faIcon{cube}\hspace{0.5em}}{}

\end{tcolorbox}

\begin{tcolorbox}
    \subsection*{Dandy \hfill \small{6 \trans{reputation_label}}}
    \tcblower
\textit{Eitelkeit}: Du bist über alle Maßen eitel und zeigst dies auch oft und gerne.


\vspace{3mm}

\begin{tabular}{@{}ll@{}}
    \textbf{Attraktivität} & +2\\
    \textbf{Charme} & +1\\
    \textbf{Kommunikation} & +1\\
    \textbf{Etikette} & +1\\
\end{tabular}

\wbif{phasesix}{
    \vspace{3mm}
\faIcon{cube}\hspace{0.5em}}{}

\end{tcolorbox}

\begin{tcolorbox}
    \subsection*{Mauerblümchen \hfill \small{1 \trans{reputation_label}}}
    \tcblower




\begin{tabular}{@{}ll@{}}
    \textbf{Attraktivität} & -1\\
    \textbf{Charme} & -1\\
    \textbf{Kommunikation} & -1\\
    \textbf{Heimlichkeit} & +2\\
\end{tabular}

\wbif{phasesix}{
    \vspace{3mm}
\faIcon{cube}\hspace{0.5em}}{}

\end{tcolorbox}

\begin{tcolorbox}
    \subsection*{Besserwisser \hfill \small{4 \trans{reputation_label}}}
    \tcblower




\begin{tabular}{@{}ll@{}}
    \textbf{Charme} & -1\\
    \textbf{Gewissenhaftigkeit} & +1\\
    \textbf{Überzeugen} & +2\\
\end{tabular}

\wbif{phasesix}{
    \vspace{3mm}
\faIcon{cube}\hspace{0.5em}}{}

\end{tcolorbox}

\begin{tcolorbox}
    \subsection*{Tratschtante \hfill \small{2 \trans{reputation_label}}}
    \tcblower




\begin{tabular}{@{}ll@{}}
    \textbf{Kommunikation} & +3\\
    \textbf{Heimlichkeit} & -3\\
\end{tabular}

\wbif{phasesix}{
    \vspace{3mm}
\faIcon{cube}\hspace{0.5em}}{}

\end{tcolorbox}

\begin{tcolorbox}
    \subsection*{Egoistisch \hfill \small{2 \trans{reputation_label}}}
    \tcblower




\begin{tabular}{@{}ll@{}}
    \textbf{Schicksalswürfel} & +1\\
    \textbf{Gewissenhaftigkeit} & -1\\
\end{tabular}

\wbif{phasesix}{
    \vspace{3mm}
\faIcon{cube}\hspace{0.5em}}{}

\end{tcolorbox}

\begin{tcolorbox}
    \subsection*{Hilfsbereit \hfill \small{4 \trans{reputation_label}}}
    \tcblower




\begin{tabular}{@{}ll@{}}
    \textbf{Charme} & +1\\
    \textbf{Kommunikation} & +1\\
\end{tabular}

\wbif{phasesix}{
    \vspace{3mm}
\faIcon{cube}\hspace{0.5em}}{}

\end{tcolorbox}

\begin{tcolorbox}
    \subsection*{Gewissenhaft \hfill \small{4 \trans{reputation_label}}}
    \tcblower




\begin{tabular}{@{}ll@{}}
    \textbf{Gewissenhaftigkeit} & +2\\
\end{tabular}

\wbif{phasesix}{
    \vspace{3mm}
\faIcon{cube}\hspace{0.5em}}{}

\end{tcolorbox}

\begin{tcolorbox}
    \subsection*{Pedantisch \hfill \small{5 \trans{reputation_label}}}
    \tcblower




\begin{tabular}{@{}ll@{}}
    \textbf{Wiederholungswürfe} & +2\\
    \textbf{Gewissenhaftigkeit} & +1\\
\end{tabular}

\wbif{phasesix}{
    \vspace{3mm}
\faIcon{cube}\hspace{0.5em}}{}

\end{tcolorbox}

\begin{tcolorbox}
    \subsection*{Umweltschützer \hfill \small{6 \trans{reputation_label}}}
    \tcblower




\begin{tabular}{@{}ll@{}}
    \textbf{Wahrnehmung} & +2\\
    \textbf{Naturkunde} & +2\\
\end{tabular}

\wbif{phasesix}{
    \vspace{3mm}
\faIcon{plane}\hspace{0.5em}\faIcon{microchip}\hspace{0.5em}}{}

\end{tcolorbox}

\begin{tcolorbox}
    \subsection*{Kosmopolit \hfill \small{6 \trans{reputation_label}}}
    \tcblower




\begin{tabular}{@{}ll@{}}
    \textbf{Bildung} & +1\\
    \textbf{Darbietung} & +1\\
    \textbf{Kommunikation} & +1\\
\end{tabular}

\wbif{phasesix}{
    \vspace{3mm}
\faIcon{umbrella}\hspace{0.5em}\faIcon{flag}\hspace{0.5em}\faIcon{plane}\hspace{0.5em}\faIcon{microchip}\hspace{0.5em}}{}

\end{tcolorbox}

\begin{tcolorbox}
    \subsection*{Genügsam \hfill \small{4 \trans{reputation_label}}}
    \tcblower




\begin{tabular}{@{}ll@{}}
    \textbf{Willenskraft} & +1\\
    \textbf{Gewissenhaftigkeit} & +1\\
\end{tabular}

\wbif{phasesix}{
    \vspace{3mm}
\faIcon{cube}\hspace{0.5em}}{}

\end{tcolorbox}

\begin{tcolorbox}
    \subsection*{Arkane Initiation \hfill \small{8 \trans{reputation_label}}}
    \tcblower




\begin{tabular}{@{}ll@{}}
    \textbf{Max Arkana} & +1\\
    \textbf{Zauberpunkte} & +5\\
    \textbf{Magiekunde} & +2\\
\end{tabular}

\wbif{phasesix}{
    \vspace{3mm}
\faIcon{magic}\hspace{0.5em}}{}

\end{tcolorbox}

\begin{tcolorbox}
    \subsection*{Blutmagie \hfill \small{5 \trans{reputation_label}}}
    \tcblower
Du darfst statt Arkana auch Wunden für das Ausführen von Zaubern verwenden.

Aufgrund der Natur der Blutmagie ist es nicht möglich, mit Zaubern, die durch Blutmagie gewirkt werden, Wunden zu heilen.


\vspace{3mm}

\begin{tabular}{@{}ll@{}}
    \textbf{Zauberpunkte} & +5\\
    \textbf{Magiekunde} & +1\\
\end{tabular}

\wbif{phasesix}{
    \vspace{3mm}
\faIcon{magic}\hspace{0.5em}}{}

\end{tcolorbox}

\begin{tcolorbox}
    \subsection*{Widerstandsfähig gegen Süchte \hfill \small{3 \trans{reputation_label}}}
    \tcblower




\begin{tabular}{@{}ll@{}}
    \textbf{Gewissenhaftigkeit} & +1\\
\end{tabular}

\wbif{phasesix}{
    \vspace{3mm}
\faIcon{cube}\hspace{0.5em}}{}

\end{tcolorbox}

\begin{tcolorbox}
    \subsection*{Introvertiert \hfill \small{2 \trans{reputation_label}}}
    \tcblower




\begin{tabular}{@{}ll@{}}
    \textbf{Schicksalswürfel} & +1\\
    \textbf{Gewissenhaftigkeit} & +1\\
    \textbf{Kommunikation} & -2\\
\end{tabular}

\wbif{phasesix}{
    \vspace{3mm}
\faIcon{cube}\hspace{0.5em}}{}

\end{tcolorbox}

\begin{tcolorbox}
    \subsection*{Reaktionär \hfill \small{-2 \trans{reputation_label}}}
    \tcblower
Der Charakter ist nicht sehr offen gegenüber Fremden, Neuem und tendiert eher zu extrem konservativen Ansichten und noch reaktionäreren Weltbildern.


\vspace{3mm}

\begin{tabular}{@{}ll@{}}
    \textbf{Charme} & -1\\
\end{tabular}

\wbif{phasesix}{
    \vspace{3mm}
\faIcon{cube}\hspace{0.5em}}{}

\end{tcolorbox}

\begin{tcolorbox}
    \subsection*{Paranoid \hfill \small{1 \trans{reputation_label}}}
    \tcblower
Der Character ist von Natur aus Paranoid, traut niemandem voll und ganz, ist immer auf der Hut.


\vspace{3mm}

\begin{tabular}{@{}ll@{}}
    \textbf{Bonuswürfel} & +1\\
    \textbf{Gewissenhaftigkeit} & -2\\
\end{tabular}

\wbif{phasesix}{
    \vspace{3mm}
\faIcon{cube}\hspace{0.5em}}{}

\end{tcolorbox}

\begin{tcolorbox}
    \subsection*{Streithahn \hfill \small{6 \trans{reputation_label}}}
    \tcblower




\begin{tabular}{@{}ll@{}}
    \textbf{Maximale Lebenspunkte} & +1\\
    \textbf{Nahkampf} & +1\\
\end{tabular}

\wbif{phasesix}{
    \vspace{3mm}
\faIcon{cube}\hspace{0.5em}}{}

\end{tcolorbox}

\begin{tcolorbox}
    \subsection*{Meisterliches Selbstbewustsein \hfill \small{40 \trans{reputation_label}}}
    \tcblower
Würfelergebnisse von 1 können einmal wiederholt werden.

\begin{pquote}{Jason Day}
When you have a lot of confidence and you feel like nobody can beat you, it’s game over for everyone else.
\end{pquote}

\vspace{3mm}


\wbif{phasesix}{
    \vspace{3mm}
\faIcon{cube}\hspace{0.5em}}{}

\end{tcolorbox}

\begin{tcolorbox}
    \subsection*{Bücherwurm \hfill \small{2 \trans{reputation_label}}}
    \tcblower




\begin{tabular}{@{}ll@{}}
    \textbf{Bildung} & +1\\
    \textbf{Ausdauer} & -1\\
    \textbf{Auffassungsgabe} & +1\\
    \textbf{Kraft} & -1\\
    \textbf{Allgemeinwissen} & +1\\
\end{tabular}

\wbif{phasesix}{
    \vspace{3mm}
\faIcon{cube}\hspace{0.5em}}{}

\end{tcolorbox}

\begin{tcolorbox}
    \subsection*{Sadist \hfill \small{3 \trans{reputation_label}}}
    \tcblower




\begin{tabular}{@{}ll@{}}
    \textbf{Charme} & +1\\
    \textbf{Empathie} & +1\\
    \textbf{Überzeugen} & +1\\
\end{tabular}

\wbif{phasesix}{
    \vspace{3mm}
\faIcon{cube}\hspace{0.5em}}{}

\end{tcolorbox}

\begin{tcolorbox}
    \subsection*{Rational \hfill \small{3 \trans{reputation_label}}}
    \tcblower




\begin{tabular}{@{}ll@{}}
    \textbf{Logik} & +1\\
\end{tabular}

\wbif{phasesix}{
    \vspace{3mm}
\faIcon{cube}\hspace{0.5em}}{}

\end{tcolorbox}

\begin{tcolorbox}
    \subsection*{Einschüchternd \hfill \small{4 \trans{reputation_label}}}
    \tcblower




\begin{tabular}{@{}ll@{}}
    \textbf{Charme} & -1\\
    \textbf{Einschüchtern} & +2\\
\end{tabular}

\wbif{phasesix}{
    \vspace{3mm}
\faIcon{cube}\hspace{0.5em}}{}

\end{tcolorbox}

\begin{tcolorbox}
    \subsection*{Ekel \hfill \small{-3 \trans{reputation_label}}}
    \tcblower
Der Charakter empfindet ausgeprägten Ekel gegen ein bestimmtes Subjekt und wird sich wenn möglich davon fern halten.


\vspace{3mm}


\wbif{phasesix}{
    \vspace{3mm}
\faIcon{cube}\hspace{0.5em}}{}

\end{tcolorbox}

\begin{tcolorbox}
    \subsection*{Gerechtigkeitsfanatiker \hfill \small{1 \trans{reputation_label}}}
    \tcblower
Der Charakter ist militanter Gerechtigkeitsfanatiker. Wenn er im Spiel eine Situation erlebt, die er als ungerecht empfindet, kann er sich kaum kontrollieren.


\vspace{3mm}

\begin{tabular}{@{}ll@{}}
    \textbf{Empathie} & +1\\
\end{tabular}

\wbif{phasesix}{
    \vspace{3mm}
\faIcon{cube}\hspace{0.5em}}{}

\end{tcolorbox}

\begin{tcolorbox}
    \subsection*{Gläubig \hfill \small{5 \trans{reputation_label}}}
    \tcblower
Der Charakter ist besonders gläubig. Alle Effekte welche von Klerikern der selben Ausrichtung an ihm ausgeführt werden haben den doppelten Effekt. Kleriker mit dieser Gabe erhalten zudem mehr Gunstpunkte, wenn sie einen zeremoniellen Gottesdienst abhalten.


\vspace{3mm}


\wbif{phasesix}{
    \vspace{3mm}
\faIcon{ankh}\hspace{0.5em}}{}

\end{tcolorbox}

\begin{tcolorbox}
    \subsection*{Habgierig \hfill \small{4 \trans{reputation_label}}}
    \tcblower
Ein habgieriger Charakter hat zuerst immer seine persönliche Bereicherung im Sinn. Das beginnt damit, dass er gerne einmal versucht, Gold oder Belohnungen seiner Reisegruppe für sich einzustreichen. Es bedeutet aber auch eine fast magische Anziehungskraft von Gold und Wertgegenständen aller Art.


\vspace{3mm}

\begin{tabular}{@{}ll@{}}
    \textbf{Logik} & +1\\
    \textbf{Täuschen und Betrügen} & +1\\
\end{tabular}

\wbif{phasesix}{
    \vspace{3mm}
\faIcon{cube}\hspace{0.5em}}{}

\end{tcolorbox}

\begin{tcolorbox}
    \subsection*{Jähzornig \hfill \small{-4 \trans{reputation_label}}}
    \tcblower
Ein jähzorniger Charakter fährt schnell aus der Haut und hat ein dünnes Fell. Bei jeder Gelegenheit die für den Charakter beleidigend wirkt kann der Spielleiter einen Wurf auf Logik verlangen. Misslingt dieser, fliegen vorraussichtlich zumindest die Fäuste.


\vspace{3mm}

\begin{tabular}{@{}ll@{}}
    \textbf{Logik} & -1\\
\end{tabular}

\wbif{phasesix}{
    \vspace{3mm}
\faIcon{cube}\hspace{0.5em}}{}

\end{tcolorbox}

\begin{tcolorbox}
    \subsection*{Konfus \hfill \small{-4 \trans{reputation_label}}}
    \tcblower
Ein Charakter mit diesem Charakterzug lässt sich leicht verwirren. Auf belebten Märkten oder in Menschenmengen kann der Spielleiter einen Wurf auf Orientierung verlangen, damit der Charakter nicht in Panik verfällt.


\vspace{3mm}

\begin{tabular}{@{}ll@{}}
    \textbf{Orientierung} & -2\\
\end{tabular}

\wbif{phasesix}{
    \vspace{3mm}
\faIcon{cube}\hspace{0.5em}}{}

\end{tcolorbox}

\begin{tcolorbox}
    \subsection*{Landei \hfill \small{-2 \trans{reputation_label}}}
    \tcblower
Der Charakter kommt vom Land. Stand und Bildung sind nicht relevant, sobald der Charakter in eine größere Siedlung (ab 1000 Einwohner) kommt ist er verwirrt. Würfe auf Orientierung in größeren Siedlungen haben einen um 1 erhöhten Mindeswurf.


\vspace{3mm}


\wbif{phasesix}{
    \vspace{3mm}
\faIcon{cube}\hspace{0.5em}}{}

\end{tcolorbox}

\begin{tcolorbox}
    \subsection*{Süchtig \hfill \small{-4 \trans{reputation_label}}}
    \tcblower
Der Character ist nach einer gewissen Substanz süchtig, je nach Ausprägung der Sucht kann ihn der Verzicht oder die Aussicht bald verzichten zu müssen in seinem Tun beeinflussen.


\vspace{3mm}


\wbif{phasesix}{
    \vspace{3mm}
\faIcon{cube}\hspace{0.5em}}{}

\end{tcolorbox}

\begin{tcolorbox}
    \subsection*{Todesbote \hfill \small{-3 \trans{reputation_label}}}
    \tcblower
Der Charakter zieht das Unheil an. Kommt er in eine neue Region oder Gesellschaft, so muss er einen W6 Werfen. Zeigt der Wurf eine 5 oder 6, so geschieht ein Unheil, Unfall oder Ähnliches.


\vspace{3mm}


\wbif{phasesix}{
    \vspace{3mm}
\faIcon{cube}\hspace{0.5em}}{}

\end{tcolorbox}

\begin{tcolorbox}
    \subsection*{Waffennarr \hfill \small{5 \trans{reputation_label}}}
    \tcblower




\begin{tabular}{@{}ll@{}}
    \textbf{Nahkampf} & +1\\
    \textbf{Schießen} & +1\\
    \textbf{Kriegsführung} & +1\\
\end{tabular}

\wbif{phasesix}{
    \vspace{3mm}
\faIcon{cube}\hspace{0.5em}}{}

\end{tcolorbox}

\begin{tcolorbox}
    \subsection*{Naiv \hfill \small{-3 \trans{reputation_label}}}
    \tcblower
\textit{Naiv}: Dein Charakter ist naiv. Er glaubt manchmal zu sehr an das Gute im Menschen. Immer, wenn der Charakter an den Aussagen oder Intentionen von NSC zweifelt, kann der Spielleiter ihn auf Logik würfeln lassen. Misslingt der Wurf, so glaubt der Charakter dem NSC.


\vspace{3mm}


\wbif{phasesix}{
    \vspace{3mm}
\faIcon{cube}\hspace{0.5em}}{}

\end{tcolorbox}

\begin{tcolorbox}
    \subsection*{Aversion \hfill \small{-2 \trans{reputation_label}}}
    \tcblower
Der Charakter hat eine Abneigung gegen ein bestimmtes Subjekt. Alle Würfe, die mit dem Subjekt interagieren haben einen um 1 erhöhten Mindestwurf.


\vspace{3mm}


\wbif{phasesix}{
    \vspace{3mm}
\faIcon{cube}\hspace{0.5em}}{}

\end{tcolorbox}

\begin{tcolorbox}
    \subsection*{Willensstark \hfill \small{3 \trans{reputation_label}}}
    \tcblower




\begin{tabular}{@{}ll@{}}
    \textbf{Willenskraft} & +2\\
\end{tabular}

\wbif{phasesix}{
    \vspace{3mm}
\faIcon{cube}\hspace{0.5em}}{}

\end{tcolorbox}

\begin{tcolorbox}
    \subsection*{Gutes Benehmen \hfill \small{4 \trans{reputation_label}}}
    \tcblower




\begin{tabular}{@{}ll@{}}
    \textbf{Etikette} & +2\\
\end{tabular}

\wbif{phasesix}{
    \vspace{3mm}
\faIcon{cube}\hspace{0.5em}}{}

\end{tcolorbox}

\begin{tcolorbox}
    \subsection*{Extrovertiert \hfill \small{5 \trans{reputation_label}}}
    \tcblower




\begin{tabular}{@{}ll@{}}
    \textbf{Auffassungsgabe} & +1\\
    \textbf{Kommunikation} & +2\\
    \textbf{Überzeugen} & +5\\
\end{tabular}

\wbif{phasesix}{
    \vspace{3mm}
\faIcon{chess-rook}\hspace{0.5em}\faIcon{dragon}\hspace{0.5em}\faIcon{ankh}\hspace{0.5em}\faIcon{magic}\hspace{0.5em}}{}

\end{tcolorbox}

\section{Talent}

Talente sind spezifische Charakterfähigkeiten. Die Schablonen beziehen sich auf eng gefasste Begabungen des Charakters. In der Regel sind sie günstig oder fügen spezielle Fähigkeiten und Regeln hinzu.

\begin{tcolorbox}
    \subsection*{Revolverheld \hfill \small{5 \trans{reputation_label}}}
    \tcblower




\begin{tabular}{@{}ll@{}}
    \textbf{Schnelligkeit} & +1\\
    \textbf{Schießen} & +2\\
\end{tabular}

\wbif{phasesix}{
    \vspace{3mm}
\faIcon{umbrella}\hspace{0.5em}\faIcon{flag}\hspace{0.5em}\faIcon{plane}\hspace{0.5em}\faIcon{microchip}\hspace{0.5em}}{}

\end{tcolorbox}

\begin{tcolorbox}
    \subsection*{Empathisch \hfill \small{5 \trans{reputation_label}}}
    \tcblower




\begin{tabular}{@{}ll@{}}
    \textbf{Empathie} & +3\\
\end{tabular}

\wbif{phasesix}{
    \vspace{3mm}
\faIcon{cube}\hspace{0.5em}}{}

\end{tcolorbox}

\begin{tcolorbox}
    \subsection*{Guter Redner \hfill \small{5 \trans{reputation_label}}}
    \tcblower




\begin{tabular}{@{}ll@{}}
    \textbf{Kommunikation} & +3\\
\end{tabular}

\wbif{phasesix}{
    \vspace{3mm}
\faIcon{cube}\hspace{0.5em}}{}

\end{tcolorbox}

\begin{tcolorbox}
    \subsection*{Athletisch \hfill \small{6 \trans{reputation_label}}}
    \tcblower




\begin{tabular}{@{}ll@{}}
    \textbf{Geschick} & +2\\
    \textbf{Ausdauer} & +2\\
\end{tabular}

\wbif{phasesix}{
    \vspace{3mm}
\faIcon{cube}\hspace{0.5em}}{}

\end{tcolorbox}

\begin{tcolorbox}
    \subsection*{Glück \hfill \small{4 \trans{reputation_label}}}
    \tcblower
Der Charakter kann zweimal pro Sitzung bis zu 3 Würfel neu werfen oder diese von einem Gefährten neu werfen lassen.


\vspace{3mm}


\wbif{phasesix}{
    \vspace{3mm}
\faIcon{cube}\hspace{0.5em}}{}

\end{tcolorbox}

\begin{tcolorbox}
    \subsection*{Führer \hfill \small{5 \trans{reputation_label}}}
    \tcblower
Der Charakter darf in jeder Kampfrunde eine Aktion an einen befreundeten Charakter weitergeben statt sie selbst zu verwenden.


\vspace{3mm}

\begin{tabular}{@{}ll@{}}
    \textbf{Bonuswürfel} & +1\\
    \textbf{Ausdauer} & +1\\
    \textbf{Auffassungsgabe} & +1\\
\end{tabular}

\wbif{phasesix}{
    \vspace{3mm}
\faIcon{cube}\hspace{0.5em}}{}

\end{tcolorbox}

\begin{tcolorbox}
    \subsection*{Guter Werfer \hfill \small{5 \trans{reputation_label}}}
    \tcblower




\begin{tabular}{@{}ll@{}}
    \textbf{Werfen} & +3\\
\end{tabular}

\wbif{phasesix}{
    \vspace{3mm}
\faIcon{cube}\hspace{0.5em}}{}

\end{tcolorbox}

\begin{tcolorbox}
    \subsection*{Synästhesie \hfill \small{5 \trans{reputation_label}}}
    \tcblower




\begin{tabular}{@{}ll@{}}
    \textbf{Auffassungsgabe} & +1\\
    \textbf{Wahrnehmung} & +2\\
\end{tabular}

\wbif{phasesix}{
    \vspace{3mm}
\faIcon{cube}\hspace{0.5em}}{}

\end{tcolorbox}

\begin{tcolorbox}
    \subsection*{Schlangenmensch \hfill \small{6 \trans{reputation_label}}}
    \tcblower




\begin{tabular}{@{}ll@{}}
    \textbf{Geschick} & +3\\
    \textbf{Resistenz} & +1\\
\end{tabular}

\wbif{phasesix}{
    \vspace{3mm}
\faIcon{cube}\hspace{0.5em}}{}

\end{tcolorbox}

\begin{tcolorbox}
    \subsection*{Starkes Immunsystem \hfill \small{5 \trans{reputation_label}}}
    \tcblower




\begin{tabular}{@{}ll@{}}
    \textbf{Resistenz} & +3\\
\end{tabular}

\wbif{phasesix}{
    \vspace{3mm}
\faIcon{cube}\hspace{0.5em}}{}

\end{tcolorbox}

\begin{tcolorbox}
    \subsection*{Photographisches Gedächtnis \hfill \small{6 \trans{reputation_label}}}
    \tcblower




\begin{tabular}{@{}ll@{}}
    \textbf{Logik} & +1\\
    \textbf{Wahrnehmung} & +1\\
    \textbf{Orientierung} & +2\\
\end{tabular}

\wbif{phasesix}{
    \vspace{3mm}
\faIcon{cube}\hspace{0.5em}}{}

\end{tcolorbox}

\begin{tcolorbox}
    \subsection*{Eiskalte Händchen \hfill \small{2 \trans{reputation_label}}}
    \tcblower




\begin{tabular}{@{}ll@{}}
    \textbf{Attraktivität} & -1\\
    \textbf{Einschüchtern} & +1\\
\end{tabular}

\wbif{phasesix}{
    \vspace{3mm}
\faIcon{cube}\hspace{0.5em}}{}

\end{tcolorbox}

\begin{tcolorbox}
    \subsection*{Kraftpaket \hfill \small{5 \trans{reputation_label}}}
    \tcblower




\begin{tabular}{@{}ll@{}}
    \textbf{Kraft} & +2\\
    \textbf{Einschüchtern} & +1\\
\end{tabular}

\wbif{phasesix}{
    \vspace{3mm}
\faIcon{cube}\hspace{0.5em}}{}

\end{tcolorbox}

\begin{tcolorbox}
    \subsection*{Medium \hfill \small{4 \trans{reputation_label}}}
    \tcblower
\textit{Visionen}: Du hast unregelmäßig Visionen. Diese können durch einen Auslöser ausgelöst werden oder rein zufällig passieren.


\vspace{3mm}

\begin{tabular}{@{}ll@{}}
    \textbf{Logik} & -1\\
    \textbf{Empathie} & +1\\
    \textbf{Wahrnehmung} & +2\\
\end{tabular}

\wbif{phasesix}{
    \vspace{3mm}
\faIcon{spider}\hspace{0.5em}}{}

\end{tcolorbox}

\begin{tcolorbox}
    \subsection*{Guter Schläfer \hfill \small{3 \trans{reputation_label}}}
    \tcblower




\begin{tabular}{@{}ll@{}}
    \textbf{Rast Mindestwurf} & -1\\
    \textbf{Gewissenhaftigkeit} & +1\\
\end{tabular}

\wbif{phasesix}{
    \vspace{3mm}
\faIcon{cube}\hspace{0.5em}}{}

\end{tcolorbox}

\begin{tcolorbox}
    \subsection*{Tierempathie \hfill \small{6 \trans{reputation_label}}}
    \tcblower
Der Mindestwurf ist bei allen Würfen, die Tiere involvieren, um 2 verringert.


\vspace{3mm}

\begin{tabular}{@{}ll@{}}
    \textbf{Empathie} & +1\\
\end{tabular}

\wbif{phasesix}{
    \vspace{3mm}
\faIcon{cube}\hspace{0.5em}}{}

\end{tcolorbox}

\begin{tcolorbox}
    \subsection*{Bauernfänger \hfill \small{8 \trans{reputation_label}}}
    \tcblower
Der Charakater kann eine Person in Sicht dazu bringen einen seiner Wiederholungswürfe auf einen gerade geworfenen Wurf anzuwenden. Hierzu muss der Charakter einen Bonus- oder Schicksalswürfel ausgeben.


\vspace{3mm}

\begin{tabular}{@{}ll@{}}
    \textbf{Wiederholungswürfe} & +1\\
\end{tabular}

\wbif{phasesix}{
    \vspace{3mm}
\faIcon{cube}\hspace{0.5em}}{}

\end{tcolorbox}

\begin{tcolorbox}
    \subsection*{Joker \hfill \small{10 \trans{reputation_label}}}
    \tcblower
Der Charakter hat die Fähigkeit von besonderen Würfen zu profitieren. Wann immer der Spieler außerhalb des Kampfes eine der folgenden Würfelfiguren würfelt tritt der angegebene Effekt in Kraft:

* \textbf{Dreierpasch} - Der Charakter erhält einen Bonuswürfel
* \textbf{Kleine Straße} - Der Charakter erhält sofort 1 Boost
* \textbf{Full House} - Der Charakter erhält einen Schicksalswürfel
* \textbf{Große Straße} - Der Charakter findet einen Gegenstand in seinem Gepäck wieder (der Spieler sucht sich einen Gegenstand aus und fügt ihn seinem Gepäck hinzu)
* \textbf{Vierpasch} - Der Charakter erhält +1 Aktion in jeder Runde des nächsten Kampfes
* \textbf{Fünferpasch} - Die Gruppe des Charakters erhält eine zusätzliche Runde vor den Gegnern im nächsten Kampf

Es gilt bei einem Wurf jeweils das höchste zu erreichende Muster. 5, 4, 3, 3, 3, 2, 1 ist also eine große Strasse, aber kein Dreierpasch.

\begin{pquote}{The Joker}
As you know, madness is like gravity...all it takes is a little push.
\end{pquote}

\vspace{3mm}


\wbif{phasesix}{
    \vspace{3mm}
\faIcon{cube}\hspace{0.5em}}{}

\end{tcolorbox}

\begin{tcolorbox}
    \subsection*{Gut ausgerüstet \hfill \small{5 \trans{reputation_label}}}
    \tcblower
Der Charakter kann einen Gegenstand in seinem Gepäck finden, auch wenn dieser sich nicht darin befindet. Hierzu verzeichnet der Spieler bei dem Charakter eine Wunde und wirft einen W6:

-  \textit{1-2}: Der gewünschte Gegenstand konnte nicht gefunden werden
-  \textit{3-5}: Ein ähnlich nützlicher Gegenstand konnte gefunden werden
- \textit{6}: Der Charakter konnte exakt den Gegenstand im Gepäck finden

\begin{pquote}{MacGyver}
With a little bit of imagination, anything is possible.
\end{pquote}

\vspace{3mm}


\wbif{phasesix}{
    \vspace{3mm}
\faIcon{cube}\hspace{0.5em}}{}

\end{tcolorbox}

\begin{tcolorbox}
    \subsection*{Geschickter Kämpfer \hfill \small{10 \trans{reputation_label}}}
    \tcblower


\begin{pquote}{Joe Lewis}
Everyone has a plan until they’ve been hit.
\end{pquote}

\vspace{3mm}

\begin{tabular}{@{}ll@{}}
    \textbf{Aktionen} & +1\\
\end{tabular}

\wbif{phasesix}{
    \vspace{3mm}
\faIcon{cube}\hspace{0.5em}}{}

\end{tcolorbox}

\begin{tcolorbox}
    \subsection*{Klarträumen \hfill \small{5 \trans{reputation_label}}}
    \tcblower


\begin{pquote}{Zhuangzi}
I dreamed I was a butterfly, flitting around in the sky; then I awoke. Now I wonder: Am I a man who dreamt of being a butterfly, or am I a butterfly dreaming that I am a man?
\end{pquote}

\vspace{3mm}

\begin{tabular}{@{}ll@{}}
    \textbf{Maximaler Stress} & +1\\
    \textbf{Willenskraft} & +2\\
\end{tabular}

\wbif{phasesix}{
    \vspace{3mm}
\faIcon{spider}\hspace{0.5em}}{}

\end{tcolorbox}

\begin{tcolorbox}
    \subsection*{Scharfschütze \hfill \small{8 \trans{reputation_label}}}
    \tcblower
Einmal pro Runde kann jede 1 eines Schiessen Wurfs neu gewürfelt werden

Benötigt einen Schiessen Wert von 5 oder mehr


\vspace{3mm}

\begin{tabular}{@{}ll@{}}
    \textbf{Orientierung} & +1\\
    \textbf{Schießen} & +1\\
    \textbf{Kriegsführung} & +1\\
\end{tabular}

\wbif{phasesix}{
    \vspace{3mm}
\faIcon{cube}\hspace{0.5em}}{}

\end{tcolorbox}

\begin{tcolorbox}
    \subsection*{Kritische Treffer \hfill \small{10 \trans{reputation_label}}}
    \tcblower
Der zu erreichende Mindestwurf für kritische Treffer ist um eins verringert. Dies gilt nur für kritische Treffer, aber nicht für megakritische Treffer.


\vspace{3mm}


\wbif{phasesix}{
    \vspace{3mm}
\faIcon{cube}\hspace{0.5em}}{}

\end{tcolorbox}

\begin{tcolorbox}
    \subsection*{Inspirierender Anführer \hfill \small{15 \trans{reputation_label}}}
    \tcblower
Als eine Aktion im Kampf kannst du auf deinen Überzeugen Wert würfeln. Eine Mitglied deiner Gruppe bekommt zusätzliche Erfolge in der Höhe deiner Erfolge oder mindestens 1


\vspace{3mm}

\begin{tabular}{@{}ll@{}}
    \textbf{Charme} & +1\\
    \textbf{Überzeugen} & +1\\
\end{tabular}

\wbif{phasesix}{
    \vspace{3mm}
\faIcon{cube}\hspace{0.5em}}{}

\end{tcolorbox}

\begin{tcolorbox}
    \subsection*{Taschenspieler Trick \hfill \small{8 \trans{reputation_label}}}
    \tcblower
Als eine Reaktion im Kampf kannst du auf deine Täuschen Fähigkeit würfeln. Ein Gegner verliert soviele Erfolge wie du Erfolge erziehlst

Benötigt eine Täuschen Fähigkeit von 3 oder höher


\vspace{3mm}


\wbif{phasesix}{
    \vspace{3mm}
\faIcon{cube}\hspace{0.5em}}{}

\end{tcolorbox}

\begin{tcolorbox}
    \subsection*{Schildtraining \hfill \small{5 \trans{reputation_label}}}
    \tcblower
Der Character hat gelernt einen Schild effizent im Kampf zu führen.

Dies erlaubt dem Charakter, Angriffe mit einem Schild auch gemäß der Regel Nahkampfangriffe parieren zu parieren. Hierzu wird wahlweise Kraft oder Geschick als Fertigkeitswert verwendet.


\vspace{3mm}


\wbif{phasesix}{
    \vspace{3mm}
\faIcon{cube}\hspace{0.5em}}{}

\end{tcolorbox}

\begin{tcolorbox}
    \subsection*{Adlerauge \hfill \small{7 \trans{reputation_label}}}
    \tcblower
Reichweiten von Fernkampfwaffen sind um 25\% erhöht.


\vspace{3mm}


\wbif{phasesix}{
    \vspace{3mm}
\faIcon{cube}\hspace{0.5em}}{}

\end{tcolorbox}

\begin{tcolorbox}
    \subsection*{Ausweichen \hfill \small{11 \trans{reputation_label}}}
    \tcblower


\begin{pquote}{Frank Herbert, Dune}
Knowing where the trap is—that's the first step in evading it.
\end{pquote}

\vspace{3mm}

\begin{tabular}{@{}ll@{}}
    \textbf{Ausweichen} & +2\\
    \textbf{Schnelligkeit} & +1\\
\end{tabular}

\wbif{phasesix}{
    \vspace{3mm}
\faIcon{cube}\hspace{0.5em}}{}

\end{tcolorbox}

\begin{tcolorbox}
    \subsection*{Magisch Begabt \hfill \small{4 \trans{reputation_label}}}
    \tcblower


\begin{pquote}{Terry Pratchett}
Scientists have calculated that the chances of something so patently absurd actually existing are millions to one. But magicians have calculated that million-to-one chances crop up nine times out of ten
\end{pquote}

\vspace{3mm}

\begin{tabular}{@{}ll@{}}
    \textbf{Max Arkana} & +2\\
\end{tabular}

\wbif{phasesix}{
    \vspace{3mm}
\faIcon{magic}\hspace{0.5em}}{}

\end{tcolorbox}

\begin{tcolorbox}
    \subsection*{Tiefe Konzentration \hfill \small{6 \trans{reputation_label}}}
    \tcblower




\begin{tabular}{@{}ll@{}}
    \textbf{Max Arkana} & +3\\
\end{tabular}

\wbif{phasesix}{
    \vspace{3mm}
\faIcon{magic}\hspace{0.5em}}{}

\end{tcolorbox}

\begin{tcolorbox}
    \subsection*{Haudegen \hfill \small{4 \trans{reputation_label}}}
    \tcblower




\begin{tabular}{@{}ll@{}}
    \textbf{Nahkampf} & +1\\
    \textbf{Schießen} & +1\\
\end{tabular}

\wbif{phasesix}{
    \vspace{3mm}
\faIcon{cube}\hspace{0.5em}}{}

\end{tcolorbox}

\begin{tcolorbox}
    \subsection*{Läufer \hfill \small{6 \trans{reputation_label}}}
    \tcblower




\begin{tabular}{@{}ll@{}}
    \textbf{Schnelligkeit} & +1\\
    \textbf{Ausdauer} & +1\\
\end{tabular}

\wbif{phasesix}{
    \vspace{3mm}
\faIcon{cube}\hspace{0.5em}}{}

\end{tcolorbox}

\begin{tcolorbox}
    \subsection*{Akrobat \hfill \small{3 \trans{reputation_label}}}
    \tcblower




\begin{tabular}{@{}ll@{}}
    \textbf{Akrobatik} & +2\\
\end{tabular}

\wbif{phasesix}{
    \vspace{3mm}
\faIcon{cube}\hspace{0.5em}}{}

\end{tcolorbox}

\begin{tcolorbox}
    \subsection*{Meisterliche Präsenz \hfill \small{40 \trans{reputation_label}}}
    \tcblower
Der Mindestwurf des Charakters ist um 1 verringert.

\begin{pquote}{Gail Sheehy}
Ah, mastery... what a profoundly satisfying feeling when one finally gets on top of a new set of skills... and then sees the light under the new door those skills can open, even as another door is closing.
\end{pquote}

\vspace{3mm}

\begin{tabular}{@{}ll@{}}
    \textbf{Mindestwurf} & -1\\
\end{tabular}

\wbif{phasesix}{
    \vspace{3mm}
\faIcon{cube}\hspace{0.5em}}{}

\end{tcolorbox}

\begin{tcolorbox}
    \subsection*{Meisterliche Zuversicht \hfill \small{40 \trans{reputation_label}}}
    \tcblower
im Sinne der Exploding Dice Regel wird jedes Mal, wenn eine 6 auf einem Würfel geworfen wurde, ein Erfolg zu den Erfolgen des Wurfs hinzugefügt.

\begin{pquote}{The Dalai Lama}
With realization of one’s own potential and self-confidence in one’s ability, one can build a better world.
\end{pquote}

\vspace{3mm}


\wbif{phasesix}{
    \vspace{3mm}
\faIcon{cube}\hspace{0.5em}}{}

\end{tcolorbox}

\begin{tcolorbox}
    \subsection*{Recherche \hfill \small{4 \trans{reputation_label}}}
    \tcblower


\begin{pquote}{Sherlock Holmes}
How often have I said to you that when you have eliminated the impossible, whatever remains, however improbable, must be the truth?
\end{pquote}

\vspace{3mm}

\begin{tabular}{@{}ll@{}}
    \textbf{Untersuchen} & +2\\
\end{tabular}

\wbif{phasesix}{
    \vspace{3mm}
\faIcon{umbrella}\hspace{0.5em}\faIcon{flag}\hspace{0.5em}\faIcon{plane}\hspace{0.5em}\faIcon{microchip}\hspace{0.5em}}{}

\end{tcolorbox}

\begin{tcolorbox}
    \subsection*{Verführungskünstler \hfill \small{10 \trans{reputation_label}}}
    \tcblower
Bei jeder Charme-Probe zur Betörung 4 Bonus sowie 2 Bonuswürfel für alle Proben auf Attraktivität.


\vspace{3mm}


\wbif{phasesix}{
    \vspace{3mm}
\faIcon{cube}\hspace{0.5em}}{}

\end{tcolorbox}

\begin{tcolorbox}
    \subsection*{Schlossknacken \hfill \small{4 \trans{reputation_label}}}
    \tcblower




\begin{tabular}{@{}ll@{}}
    \textbf{Schlossknacken} & +2\\
\end{tabular}

\wbif{phasesix}{
    \vspace{3mm}
\faIcon{cube}\hspace{0.5em}}{}

\end{tcolorbox}

\begin{tcolorbox}
    \subsection*{Erste Hilfe \hfill \small{4 \trans{reputation_label}}}
    \tcblower


\begin{pquote}{Henry Dunant}
Our real enemy is not our neighboring country; it's hunger, cold, poverty, ignorance, superstition and prejudice.
\end{pquote}

\vspace{3mm}

\begin{tabular}{@{}ll@{}}
    \textbf{Erste Hilfe} & +2\\
\end{tabular}

\wbif{phasesix}{
    \vspace{3mm}
\faIcon{cube}\hspace{0.5em}}{}

\end{tcolorbox}

\begin{tcolorbox}
    \subsection*{Preschen \hfill \small{6 \trans{reputation_label}}}
    \tcblower
Immer, wenn der Charakter im Kampf die Laufen Aktion nutzt kann er einen Akrobatikwurf machen, um weitere Meter voran zu kommen. Gelingt dieser Wurf, so darf er sich gemäß den Erfolgen weiter bewegen als seine Laufen Reichweite. 

Zeigt der Wurf keinen Erfolg, so stolpert der Charakter, und gilt als Liegend. Er muss eine Aktion aufwenden, um wieder auf die Beine oder in die Hocke zu kommen.


\vspace{3mm}

\begin{tabular}{@{}ll@{}}
    \textbf{Schnelligkeit} & +1\\
\end{tabular}

\wbif{phasesix}{
    \vspace{3mm}
\faIcon{cube}\hspace{0.5em}}{}

\end{tcolorbox}

\begin{tcolorbox}
    \subsection*{Trainierter Schwertarm \hfill \small{5 \trans{reputation_label}}}
    \tcblower


\begin{pquote}{None}
"Nur hartes Training und unbändige Disziplin sorgen für Tod und Verderben in einem flüssigen Streich"
\end{pquote}

\vspace{3mm}

\begin{tabular}{@{}ll@{}}
    \textbf{Kraft} & +2\\
    \textbf{Nahkampf} & +1\\
\end{tabular}

\wbif{phasesix}{
    \vspace{3mm}
\faIcon{cube}\hspace{0.5em}}{}

\end{tcolorbox}

\begin{tcolorbox}
    \subsection*{Geschickt \hfill \small{4 \trans{reputation_label}}}
    \tcblower




\begin{tabular}{@{}ll@{}}
    \textbf{Geschick} & +2\\
\end{tabular}

\wbif{phasesix}{
    \vspace{3mm}
\faIcon{cube}\hspace{0.5em}}{}

\end{tcolorbox}

\begin{tcolorbox}
    \subsection*{Gewissenhaft \hfill \small{4 \trans{reputation_label}}}
    \tcblower




\begin{tabular}{@{}ll@{}}
    \textbf{Gewissenhaftigkeit} & +2\\
\end{tabular}

\wbif{phasesix}{
    \vspace{3mm}
\faIcon{cube}\hspace{0.5em}}{}

\end{tcolorbox}

\begin{tcolorbox}
    \subsection*{Guter Schütze \hfill \small{3 \trans{reputation_label}}}
    \tcblower




\begin{tabular}{@{}ll@{}}
    \textbf{Schießen} & +1\\
\end{tabular}

\wbif{phasesix}{
    \vspace{3mm}
\faIcon{cube}\hspace{0.5em}}{}

\end{tcolorbox}

\begin{tcolorbox}
    \subsection*{Eins mit der Magie \hfill \small{10 \trans{reputation_label}}}
    \tcblower




\begin{tabular}{@{}ll@{}}
    \textbf{Zauberpunkte} & +15\\
\end{tabular}

\wbif{phasesix}{
    \vspace{3mm}
\faIcon{magic}\hspace{0.5em}}{}

\end{tcolorbox}

\begin{tcolorbox}
    \subsection*{Arkane Erholung \hfill \small{10 \trans{reputation_label}}}
    \tcblower
Einmal am Tag kannst du außerhalb einer Rast deine Verbindung zur Magie erneuern und dein \textbf{Arkana} auffüllen. Dazu musst du etwa eine Stunde Ruhe finden, um dich auf den Fluss der Magie zu konzentrieren.

Würfel auf deinen \textbf{Magiekunde} Wert und fülle mit der Hälfte der Erfolge deine Arkana auf.


\vspace{3mm}


\wbif{phasesix}{
    \vspace{3mm}
\faIcon{magic}\hspace{0.5em}}{}

\end{tcolorbox}

\begin{tcolorbox}
    \subsection*{Arkane Meisterschaft \hfill \small{40 \trans{reputation_label}}}
    \tcblower




\begin{tabular}{@{}ll@{}}
    \textbf{Max Arkana} & +4\\
    \textbf{Zauberpunkte} & +50\\
\end{tabular}

\wbif{phasesix}{
    \vspace{3mm}
\faIcon{magic}\hspace{0.5em}}{}

\end{tcolorbox}

\begin{tcolorbox}
    \subsection*{Medizin \hfill \small{5 \trans{reputation_label}}}
    \tcblower




\begin{tabular}{@{}ll@{}}
    \textbf{Erste Hilfe} & +1\\
    \textbf{Medizin} & +2\\
\end{tabular}

\wbif{phasesix}{
    \vspace{3mm}
\faIcon{cube}\hspace{0.5em}}{}

\end{tcolorbox}

\begin{tcolorbox}
    \subsection*{Besondere Wachsamkeit \hfill \small{10 \trans{reputation_label}}}
    \tcblower
Zu Beginn eines Kampfes erhält der Charakter eine Aktion, welche allerdings nur zum Reagieren verwendet werden kann. Sobald der Charakter in der ersten Kampfrunde an der Reihe ist überschreiben seine Aktionen diese zusätzliche Aktion.


\vspace{3mm}


\wbif{phasesix}{
    \vspace{3mm}
\faIcon{cube}\hspace{0.5em}}{}

\end{tcolorbox}

\begin{tcolorbox}
    \subsection*{Verhör \hfill \small{6 \trans{reputation_label}}}
    \tcblower




\begin{tabular}{@{}ll@{}}
    \textbf{Einschüchtern} & +1\\
    \textbf{Empathie} & +1\\
    \textbf{Verhören} & +2\\
\end{tabular}

\wbif{phasesix}{
    \vspace{3mm}
\faIcon{cube}\hspace{0.5em}}{}

\end{tcolorbox}

\begin{tcolorbox}
    \subsection*{Tarnen und Verstecken \hfill \small{4 \trans{reputation_label}}}
    \tcblower




\begin{tabular}{@{}ll@{}}
    \textbf{Heimlichkeit} & +2\\
\end{tabular}

\wbif{phasesix}{
    \vspace{3mm}
\faIcon{cube}\hspace{0.5em}}{}

\end{tcolorbox}

\begin{tcolorbox}
    \subsection*{Zauberei \hfill \small{7 \trans{reputation_label}}}
    \tcblower




\begin{tabular}{@{}ll@{}}
    \textbf{Zauberpunkte} & +5\\
    \textbf{Zauberei} & True\\
\end{tabular}

\wbif{phasesix}{
    \vspace{3mm}
\faIcon{magic}\hspace{0.5em}}{}

\end{tcolorbox}

\begin{tcolorbox}
    \subsection*{Weiße Magie \hfill \small{7 \trans{reputation_label}}}
    \tcblower




\begin{tabular}{@{}ll@{}}
    \textbf{Zauberpunkte} & +7\\
    \textbf{Weisse Magie} & True\\
\end{tabular}

\wbif{phasesix}{
    \vspace{3mm}
\faIcon{magic}\hspace{0.5em}}{}

\end{tcolorbox}

\begin{tcolorbox}
    \subsection*{Schwarze Magie \hfill \small{7 \trans{reputation_label}}}
    \tcblower




\begin{tabular}{@{}ll@{}}
    \textbf{Zauberpunkte} & +5\\
    \textbf{Schwarze Magie} & True\\
\end{tabular}

\wbif{phasesix}{
    \vspace{3mm}
\faIcon{magic}\hspace{0.5em}}{}

\end{tcolorbox}

\begin{tcolorbox}
    \subsection*{Elementarmagie \hfill \small{7 \trans{reputation_label}}}
    \tcblower




\begin{tabular}{@{}ll@{}}
    \textbf{Zauberpunkte} & +5\\
    \textbf{Elementarmagie} & True\\
\end{tabular}

\wbif{phasesix}{
    \vspace{3mm}
\faIcon{magic}\hspace{0.5em}}{}

\end{tcolorbox}

\begin{tcolorbox}
    \subsection*{Schamanismus \hfill \small{7 \trans{reputation_label}}}
    \tcblower




\begin{tabular}{@{}ll@{}}
    \textbf{Zauberpunkte} & +7\\
    \textbf{Schamanismus} & True\\
\end{tabular}

\wbif{phasesix}{
    \vspace{3mm}
\faIcon{magic}\hspace{0.5em}}{}

\end{tcolorbox}

\begin{tcolorbox}
    \subsection*{Sanguine Magie \hfill \small{7 \trans{reputation_label}}}
    \tcblower




\begin{tabular}{@{}ll@{}}
    \textbf{Zauberpunkte} & +5\\
    \textbf{Sanguine Magie} & True\\
\end{tabular}

\wbif{phasesix}{
    \vspace{3mm}
\faIcon{magic}\hspace{0.5em}}{}

\end{tcolorbox}

\begin{tcolorbox}
    \subsection*{Nekrologie \hfill \small{7 \trans{reputation_label}}}
    \tcblower




\begin{tabular}{@{}ll@{}}
    \textbf{Zauberpunkte} & +5\\
    \textbf{Nekrologie} & True\\
\end{tabular}

\wbif{phasesix}{
    \vspace{3mm}
\faIcon{magic}\hspace{0.5em}}{}

\end{tcolorbox}

\begin{tcolorbox}
    \subsection*{Mystizismus \hfill \small{7 \trans{reputation_label}}}
    \tcblower




\begin{tabular}{@{}ll@{}}
    \textbf{Zauberpunkte} & +5\\
    \textbf{Mystizismus} & True\\
\end{tabular}

\wbif{phasesix}{
    \vspace{3mm}
\faIcon{magic}\hspace{0.5em}}{}

\end{tcolorbox}

\begin{tcolorbox}
    \subsection*{Hermetik \hfill \small{7 \trans{reputation_label}}}
    \tcblower




\begin{tabular}{@{}ll@{}}
    \textbf{Zauberpunkte} & +5\\
    \textbf{Hermetik} & True\\
\end{tabular}

\wbif{phasesix}{
    \vspace{3mm}
\faIcon{magic}\hspace{0.5em}}{}

\end{tcolorbox}

\begin{tcolorbox}
    \subsection*{Geisterbeschwörungen \hfill \small{7 \trans{reputation_label}}}
    \tcblower




\begin{tabular}{@{}ll@{}}
    \textbf{Zauberpunkte} & +5\\
    \textbf{Geisterbeschwörungen} & True\\
\end{tabular}

\wbif{phasesix}{
    \vspace{3mm}
\faIcon{magic}\hspace{0.5em}}{}

\end{tcolorbox}

\begin{tcolorbox}
    \subsection*{Dämonologie \hfill \small{7 \trans{reputation_label}}}
    \tcblower




\begin{tabular}{@{}ll@{}}
    \textbf{Zauberpunkte} & +5\\
    \textbf{Dämonologie} & True\\
\end{tabular}

\wbif{phasesix}{
    \vspace{3mm}
\faIcon{magic}\hspace{0.5em}}{}

\end{tcolorbox}

\begin{tcolorbox}
    \subsection*{Astrale Magie \hfill \small{7 \trans{reputation_label}}}
    \tcblower




\begin{tabular}{@{}ll@{}}
    \textbf{Zauberpunkte} & +5\\
    \textbf{Astrale Magie} & True\\
\end{tabular}

\wbif{phasesix}{
    \vspace{3mm}
\faIcon{magic}\hspace{0.5em}}{}

\end{tcolorbox}

\begin{tcolorbox}
    \subsection*{Chimärologie \hfill \small{7 \trans{reputation_label}}}
    \tcblower




\begin{tabular}{@{}ll@{}}
    \textbf{Zauberpunkte} & +5\\
    \textbf{Chimärologie} & True\\
\end{tabular}

\wbif{phasesix}{
    \vspace{3mm}
\faIcon{magic}\hspace{0.5em}}{}

\end{tcolorbox}

\begin{tcolorbox}
    \subsection*{Guter Looter \hfill \small{6 \trans{reputation_label}}}
    \tcblower
Der Charakter kann bei einem Untersuchenwurf durch einen kritischen Erfolg besonders interessante  Dinge entdecken.


\vspace{3mm}


\wbif{phasesix}{
    \vspace{3mm}
\faIcon{cube}\hspace{0.5em}}{}

\end{tcolorbox}

\begin{tcolorbox}
    \subsection*{Arkanes Naturtalent \hfill \small{6 \trans{reputation_label}}}
    \tcblower




\begin{tabular}{@{}ll@{}}
    \textbf{Zauberpunkte} & +5\\
    \textbf{Max Arkana} & +2\\
\end{tabular}

\wbif{phasesix}{
    \vspace{3mm}
\faIcon{magic}\hspace{0.5em}}{}

\end{tcolorbox}

\begin{tcolorbox}
    \subsection*{Schmerz ignorieren \hfill \small{7 \trans{reputation_label}}}
    \tcblower
Einmal pro Kampf kann auf Resistenz geworfen werden. Bei Erfolg wird der Schaden einer Angriffsquelle in dieser Aktion komplett verhindert. Dieser Wurf benötigt keine Aktion.


\vspace{3mm}


\wbif{phasesix}{
    \vspace{3mm}
\faIcon{cube}\hspace{0.5em}}{}

\end{tcolorbox}

\begin{tcolorbox}
    \subsection*{Verfluchungen \hfill \small{7 \trans{reputation_label}}}
    \tcblower




\begin{tabular}{@{}ll@{}}
    \textbf{Zauberpunkte} & +7\\
    \textbf{Verfluchungen} & True\\
\end{tabular}

\wbif{phasesix}{
    \vspace{3mm}
\faIcon{magic}\hspace{0.5em}}{}

\end{tcolorbox}

\begin{tcolorbox}
    \subsection*{Intelligent \hfill \small{4 \trans{reputation_label}}}
    \tcblower


\begin{pquote}{Albert Einstein}
Phantasie ist wichtiger als Wissen, denn Wissen ist begrenzt.
\end{pquote}

\vspace{3mm}

\begin{tabular}{@{}ll@{}}
    \textbf{Logik} & +1\\
    \textbf{Bildung} & +1\\
\end{tabular}

\wbif{phasesix}{
    \vspace{3mm}
\faIcon{cube}\hspace{0.5em}}{}

\end{tcolorbox}

\begin{tcolorbox}
    \subsection*{Agil \hfill \small{3 \trans{reputation_label}}}
    \tcblower




\begin{tabular}{@{}ll@{}}
    \textbf{Schnelligkeit} & +1\\
\end{tabular}

\wbif{phasesix}{
    \vspace{3mm}
\faIcon{cube}\hspace{0.5em}}{}

\end{tcolorbox}

\begin{tcolorbox}
    \subsection*{Schildmeisterschaft \hfill \small{10 \trans{reputation_label}}}
    \tcblower
Der Charakter ist ein Meister im Umgang mit Schilden. Der Schild kann für einen Schildblock mit einer statt zwei Aktionen vorbereitet werden. Außerdem kann der Schildblock auch als Reaktion ausgeführt werden.


\vspace{3mm}


\wbif{phasesix}{
    \vspace{3mm}
\faIcon{cube}\hspace{0.5em}}{}

\end{tcolorbox}

\begin{tcolorbox}
    \subsection*{Klingentanz \hfill \small{15 \trans{reputation_label}}}
    \tcblower
Der Charakter ist geübt im beidhändigen Kampf mit zwei Waffen. Der Mindestwurf für den Angriff mit der Sekundärwaffe ist nicht mehr um eins erhöht.


\vspace{3mm}


\wbif{phasesix}{
    \vspace{3mm}
\faIcon{cube}\hspace{0.5em}}{}

\end{tcolorbox}

\begin{tcolorbox}
    \subsection*{Gelegenheitsangriff \hfill \small{10 \trans{reputation_label}}}
    \tcblower
Der Charakter kann einmal pro Kampfrunde als Reaktion einen Angriff gegen einen Gegner ausführen, wenn dieser durch Bewegung den Wirkungsbereich (Reichweite) seiner Nahkampfwaffe verlässt. Dafür wird keine Aktion verwendet.


\vspace{3mm}


\wbif{phasesix}{
    \vspace{3mm}
\faIcon{cube}\hspace{0.5em}}{}

\end{tcolorbox}

\begin{tcolorbox}
    \subsection*{Ausweichen \hfill \small{7 \trans{reputation_label}}}
    \tcblower




\begin{tabular}{@{}ll@{}}
    \textbf{Ausweichen} & +2\\
\end{tabular}

\wbif{phasesix}{
    \vspace{3mm}
\faIcon{cube}\hspace{0.5em}}{}

\end{tcolorbox}

\begin{tcolorbox}
    \subsection*{Schnell \hfill \small{4 \trans{reputation_label}}}
    \tcblower




\begin{tabular}{@{}ll@{}}
    \textbf{Schnelligkeit} & +2\\
\end{tabular}

\wbif{phasesix}{
    \vspace{3mm}
\faIcon{cube}\hspace{0.5em}}{}

\end{tcolorbox}

\begin{tcolorbox}
    \subsection*{Vollkommene Untergebenheit \hfill \small{20 \trans{reputation_label}}}
    \tcblower
Jegliche erhaltene Gunst wird verdoppelt.


\vspace{3mm}


\wbif{phasesix}{
    \vspace{3mm}
\faIcon{ankh}\hspace{0.5em}}{}

\end{tcolorbox}

\begin{tcolorbox}
    \subsection*{Widerstand \hfill \small{20 \trans{reputation_label}}}
    \tcblower
Der Charakter ist in der Lage, Verletzungen natürlich zu widerstehen. Für jede erfolgte Wunde wird ein W6 geworfen. Der Mindestwurf hierfür beträgt 5 + die Anzahl der verursachten Wunden. Jeder Erfolg verhindert eine Wunde.


\vspace{3mm}


\wbif{phasesix}{
    \vspace{3mm}
\faIcon{cube}\hspace{0.5em}}{}

\end{tcolorbox}

\begin{tcolorbox}
    \subsection*{Bewegen \hfill \small{5 \trans{reputation_label}}}
    \tcblower




\begin{tabular}{@{}ll@{}}
    \textbf{Akrobatik} & +2\\
    \textbf{Heimlichkeit} & +2\\
\end{tabular}

\wbif{phasesix}{
    \vspace{3mm}
\faIcon{chess-rook}\hspace{0.5em}\faIcon{dragon}\hspace{0.5em}\faIcon{ankh}\hspace{0.5em}\faIcon{magic}\hspace{0.5em}}{}

\end{tcolorbox}

\begin{tcolorbox}
    \subsection*{Nahkampf \hfill \small{3 \trans{reputation_label}}}
    \tcblower




\begin{tabular}{@{}ll@{}}
    \textbf{Nahkampf} & +3\\
\end{tabular}

\wbif{phasesix}{
    \vspace{3mm}
\faIcon{chess-rook}\hspace{0.5em}\faIcon{dragon}\hspace{0.5em}\faIcon{ankh}\hspace{0.5em}\faIcon{magic}\hspace{0.5em}}{}

\end{tcolorbox}

\section{Lebensumstände}

Diese Schablonen beschreiben die Lebensumstände des Charakters.

\begin{tcolorbox}
    \subsection*{Waise \hfill \small{4 \trans{reputation_label}}}
    \tcblower




\begin{tabular}{@{}ll@{}}
    \textbf{Schicksalswürfel} & +1\\
    \textbf{Resistenz} & +1\\
\end{tabular}

\wbif{phasesix}{
    \vspace{3mm}
\faIcon{cube}\hspace{0.5em}}{}

\end{tcolorbox}

\begin{tcolorbox}
    \subsection*{Gesucht \hfill \small{4 \trans{reputation_label}}}
    \tcblower
\textit{Rivale}: Du hast einen Rivalen. Irgendjemand konkurriert mit dir, sei es auf einer beruflichen, professionellen oder tödlichen Ebene.


\vspace{3mm}

\begin{tabular}{@{}ll@{}}
    \textbf{Schicksalswürfel} & +1\\
    \textbf{Gewissenhaftigkeit} & +1\\
\end{tabular}

\wbif{phasesix}{
    \vspace{3mm}
\faIcon{cube}\hspace{0.5em}}{}

\end{tcolorbox}

\begin{tcolorbox}
    \subsection*{Alleinerziehend \hfill \small{4 \trans{reputation_label}}}
    \tcblower




\begin{tabular}{@{}ll@{}}
    \textbf{Willenskraft} & +1\\
    \textbf{Tapferkeit} & +1\\
\end{tabular}

\wbif{phasesix}{
    \vspace{3mm}
\faIcon{cube}\hspace{0.5em}}{}

\end{tcolorbox}

\begin{tcolorbox}
    \subsection*{Verlust eines Familienmitgliedes \hfill \small{5 \trans{reputation_label}}}
    \tcblower




\begin{tabular}{@{}ll@{}}
    \textbf{Schicksalswürfel} & +1\\
    \textbf{Tapferkeit} & +2\\
\end{tabular}

\wbif{phasesix}{
    \vspace{3mm}
\faIcon{cube}\hspace{0.5em}}{}

\end{tcolorbox}

\begin{tcolorbox}
    \subsection*{Adelig \hfill \small{9 \trans{reputation_label}}}
    \tcblower




\begin{tabular}{@{}ll@{}}
    \textbf{Maximale Lebenspunkte} & +1\\
    \textbf{Kommunikation} & +2\\
    \textbf{Nahkampf} & +1\\
    \textbf{Täuschen und Betrügen} & +1\\
    \textbf{Lesen/Schreiben} & +2\\
\end{tabular}

\wbif{phasesix}{
    \vspace{3mm}
\faIcon{chess-rook}\hspace{0.5em}}{}

\end{tcolorbox}

\begin{tcolorbox}
    \subsection*{Verlust eines Körperteils \hfill \small{1 \trans{reputation_label}}}
    \tcblower




\begin{tabular}{@{}ll@{}}
    \textbf{Schicksalswürfel} & +1\\
    \textbf{Attraktivität} & -1\\
    \textbf{Geschick} & -1\\
\end{tabular}

\wbif{phasesix}{
    \vspace{3mm}
\faIcon{cube}\hspace{0.5em}}{}

\end{tcolorbox}

\begin{tcolorbox}
    \subsection*{Guru \hfill \small{7 \trans{reputation_label}}}
    \tcblower




\begin{tabular}{@{}ll@{}}
    \textbf{Charme} & +1\\
    \textbf{Attraktivität} & +2\\
    \textbf{Kommunikation} & +2\\
\end{tabular}

\wbif{phasesix}{
    \vspace{3mm}
\faIcon{cube}\hspace{0.5em}}{}

\end{tcolorbox}

\begin{tcolorbox}
    \subsection*{Entstellt \hfill \small{4 \trans{reputation_label}}}
    \tcblower




\begin{tabular}{@{}ll@{}}
    \textbf{Charme} & -1\\
    \textbf{Attraktivität} & -1\\
    \textbf{Einschüchtern} & +2\\
    \textbf{Tapferkeit} & +2\\
\end{tabular}

\wbif{phasesix}{
    \vspace{3mm}
\faIcon{cube}\hspace{0.5em}}{}

\end{tcolorbox}

\begin{tcolorbox}
    \subsection*{Unglücklich verliebt \hfill \small{3 \trans{reputation_label}}}
    \tcblower




\begin{tabular}{@{}ll@{}}
    \textbf{Willenskraft} & +1\\
\end{tabular}

\wbif{phasesix}{
    \vspace{3mm}
\faIcon{cube}\hspace{0.5em}}{}

\end{tcolorbox}

\begin{tcolorbox}
    \subsection*{Leibeigener \hfill \small{5 \trans{reputation_label}}}
    \tcblower
\textit{Obrigkeitshörig}: Du befolgst jeden Befehl deines Vorgesetzten, ohne weiter darüber nachzudenken.


\vspace{3mm}

\begin{tabular}{@{}ll@{}}
    \textbf{Auffassungsgabe} & +2\\
    \textbf{Heimlichkeit} & +1\\
\end{tabular}

\wbif{phasesix}{
    \vspace{3mm}
\faIcon{chess-rook}\hspace{0.5em}\faIcon{umbrella}\hspace{0.5em}}{}

\end{tcolorbox}

\begin{tcolorbox}
    \subsection*{Klosterleben \hfill \small{5 \trans{reputation_label}}}
    \tcblower




\begin{tabular}{@{}ll@{}}
    \textbf{Gewissenhaftigkeit} & +1\\
    \textbf{Naturkunde} & +1\\
    \textbf{Kommunikation} & -1\\
    \textbf{Religion} & +2\\
\end{tabular}

\wbif{phasesix}{
    \vspace{3mm}
\faIcon{cube}\hspace{0.5em}}{}

\end{tcolorbox}

\begin{tcolorbox}
    \subsection*{Eremit \hfill \small{2 \trans{reputation_label}}}
    \tcblower




\begin{tabular}{@{}ll@{}}
    \textbf{Charme} & -1\\
    \textbf{Orientierung} & +1\\
    \textbf{Wahrnehmung} & +1\\
    \textbf{Kommunikation} & -1\\
\end{tabular}

\wbif{phasesix}{
    \vspace{3mm}
\faIcon{cube}\hspace{0.5em}}{}

\end{tcolorbox}

\begin{tcolorbox}
    \subsection*{Obdachlos \hfill \small{5 \trans{reputation_label}}}
    \tcblower
\textit{Begleiter}: Du darfst einen tierischen Begleiter auswählen, welcher als Vertrauter gilt und dich auf Schritt und Tritt begleitet.


\vspace{3mm}

\begin{tabular}{@{}ll@{}}
    \textbf{Attraktivität} & -1\\
    \textbf{Resistenz} & +3\\
    \textbf{Nahkampf} & +1\\
\end{tabular}

\wbif{phasesix}{
    \vspace{3mm}
\faIcon{cube}\hspace{0.5em}}{}

\end{tcolorbox}

\begin{tcolorbox}
    \subsection*{Witwer \hfill \small{5 \trans{reputation_label}}}
    \tcblower




\begin{tabular}{@{}ll@{}}
    \textbf{Bonuswürfel} & +1\\
    \textbf{Tapferkeit} & +2\\
\end{tabular}

\wbif{phasesix}{
    \vspace{3mm}
\faIcon{cube}\hspace{0.5em}}{}

\end{tcolorbox}

\begin{tcolorbox}
    \subsection*{Schmiss \hfill \small{-2 \trans{reputation_label}}}
    \tcblower




\begin{tabular}{@{}ll@{}}
    \textbf{Attraktivität} & -1\\
\end{tabular}

\wbif{phasesix}{
    \vspace{3mm}
\faIcon{umbrella}\hspace{0.5em}\faIcon{flag}\hspace{0.5em}\faIcon{plane}\hspace{0.5em}\faIcon{microchip}\hspace{0.5em}}{}

\end{tcolorbox}

\begin{tcolorbox}
    \subsection*{Wanderjahre \hfill \small{7 \trans{reputation_label}}}
    \tcblower




\begin{tabular}{@{}ll@{}}
    \textbf{Bildung} & +1\\
    \textbf{Orientierung} & +1\\
    \textbf{Tapferkeit} & +2\\
    \textbf{Mechanik} & +1\\
\end{tabular}

\wbif{phasesix}{
    \vspace{3mm}
\faIcon{chess-rook}\hspace{0.5em}\faIcon{umbrella}\hspace{0.5em}}{}

\end{tcolorbox}

\begin{tcolorbox}
    \subsection*{Vampir \hfill \small{4 \trans{reputation_label}}}
    \tcblower




\begin{tabular}{@{}ll@{}}
    \textbf{Schicksalswürfel} & +2\\
    \textbf{Attraktivität} & -1\\
    \textbf{Resistenz} & +1\\
\end{tabular}

\wbif{phasesix}{
    \vspace{3mm}
\faIcon{spider}\hspace{0.5em}}{}

\end{tcolorbox}

\begin{tcolorbox}
    \subsection*{Magische Begegnung \hfill \small{5 \trans{reputation_label}}}
    \tcblower




\begin{tabular}{@{}ll@{}}
    \textbf{Zauberpunkte} & +5\\
\end{tabular}

\wbif{phasesix}{
    \vspace{3mm}
\faIcon{magic}\hspace{0.5em}}{}

\end{tcolorbox}

\begin{tcolorbox}
    \subsection*{Borkenhaut \hfill \small{3 \trans{reputation_label}}}
    \tcblower




\begin{tabular}{@{}ll@{}}
    \textbf{Schutz} & +1\\
    \textbf{Schnelligkeit} & -1\\
\end{tabular}

\wbif{phasesix}{
    \vspace{3mm}
\faIcon{magic}\hspace{0.5em}}{}

\end{tcolorbox}

\begin{tcolorbox}
    \subsection*{Vampir \hfill \small{8 \trans{reputation_label}}}
    \tcblower
\textit{Vampir}: Du bist ein Untoter. Dieser Umstand sorgt dafür, dass du zu einem Haufen Staub zerfällst wenn du stirbst. Ein Tropfen Blut kann dich dann immer wieder beleben. Der Aufenthalt in direktem Sonnenlicht lässt dich pro Stunde eine direkte Wunde hinnehmen. Wenn du eine Gottheit hast, kannst du die Handlungen eines Priesters ausführen und seine Gunst erbitten.


\vspace{3mm}

\begin{tabular}{@{}ll@{}}
    \textbf{Maximale Lebenspunkte} & +2\\
    \textbf{Altertümer} & +2\\
\end{tabular}

\wbif{phasesix}{
    \vspace{3mm}
\faIcon{magic}\hspace{0.5em}}{}

\end{tcolorbox}

\begin{tcolorbox}
    \subsection*{Gesegnet \hfill \small{60 \trans{reputation_label}}}
    \tcblower




\begin{tabular}{@{}ll@{}}
    \textbf{Mindestwurf} & -1\\
    \textbf{Schicksalswürfel} & +2\\
\end{tabular}

\wbif{phasesix}{
    \vspace{3mm}
\faIcon{ankh}\hspace{0.5em}}{}

\end{tcolorbox}

\end{multicols}